\documentclass[twoside]{article}
\usepackage[utf8]{inputenc}
\usepackage[english]{babel}
\usepackage{geometry}
\usepackage{multicol}
\usepackage{minted}
\usepackage{python}
\usepackage{fontspec}
\usepackage[hidelinks]{hyperref}
\usepackage{fancyhdr}
\usepackage{listings}
\usepackage{pdfpages}
\usepackage{needspace}
\usepackage{sectsty}
\usepackage{array}
\usepackage{multirow}
\usepackage{longtable}
\usepackage{xcolor}
\usepackage{afterpage}
\usepackage{amssymb}
\usepackage{amsmath}
\usepackage{graphicx}
\usepackage{svg}
\usepackage[inline]{enumitem}

\graphicspath{ {./media/} }

\definecolor{prussianblue}{rgb}{0.0, 0.19, 0.33}
\definecolor{indigo(dye)}{rgb}{0.0, 0.25, 0.42}
\definecolor{lapislazuli}{rgb}{0.15, 0.38, 0.61}
\definecolor{mediumelectricblue}{rgb}{0.01, 0.31, 0.59}
\definecolor{smalt(darkpowderblue)}{rgb}{0.0, 0.2, 0.6}
\definecolor{yaleblue}{rgb}{0.06, 0.3, 0.57}
\definecolor{skobeloff}{rgb}{0.0, 0.48, 0.45}
\definecolor{pinegreen}{rgb}{0.0, 0.47, 0.44}
\definecolor{coolgray}{RGB}{49, 63, 71}
\definecolor{coolpink}{RGB}{255,116,150}

\setmainfont[
    Path=fonts/,
    UprightFont = nexalight.otf,
    BoldFont = nexabold.otf
]{nexalight}

\setsansfont[
    Path=fonts/,
    UprightFont = nexalight.otf,
    BoldFont = nexabold.otf
]{nexalight}


\geometry{letterpaper, landscape, left=0.5cm, right=0.5cm, top=1.8cm, bottom=1cm}

\sectionfont{\Huge\bfseries\sffamily}

\setminted{
    style=manni,
    breaklines=true,
    breakanywhere=true,
    tabsize=4
}

\setlength{\color{lapislazuli}\headsep}{0.5cm}
\setlength{\columnsep}{0.5cm}
\setlength{\columnseprule}{0.1cm}
\renewcommand{\columnseprulecolor}{\color{lapislazuli}}


\pagenumbering{arabic}
\fancyhead{}
\fancyfoot{}
\fancyhead[LO,RE]{\textcolor{coolgray}{acuantoelmamutdeapeso}}
\fancyhead[LE,RO]{\textcolor{coolgray}{\leftmark}}
\fancyfoot[LE,RO]{\textcolor{coolgray}{\textbf{\thepage}}}

\renewcommand{\headrulewidth}{0.05cm}
\renewcommand{\footrulewidth}{0.01cm}
\newcommand\YUGE{\fontsize{35}{50}\selectfont}

\setlength{\parindent}{0em}
\setlength{\tabcolsep}{10pt} % Default value: 6pt
\renewcommand{\arraystretch}{1.5} % Default value: 1

\begin{document}
\pagestyle{fancy}
\newpage

\selectfont
\color{coolgray}
\begin{multicols*}{2}
	\tableofcontents
\end{multicols*}


\begin{multicols*}{2}
	\needspace{5\baselineskip}
	{
		\phantomsection
		\section*{Template}
		\markboth{TEMPLATE}{}
		\addcontentsline{toc}{section}{Template}
	}
	{
		\phantomsection
		\subsection*{Template}
		\addcontentsline{toc}{subsection}{Template}
	}
	\begin{minted}{cpp}
#include <bits/stdc++.h>
#include <ext/pb_ds/assoc_container.hpp>
#include <ext/pb_ds/tree_policy.hpp>
using namespace std;
using namespace __gnu_pbds;

#define all(v) v.begin(), v.end()
#define rall(v) v.rbegin(), v.rend()
#define pb push_back
#define fi first
#define se second
#define sz(a) int(a.size())
#define endl "\n"
using ll = long long;
using vi = vector<int>;
using vll = vector<ll>;
using pii = pair<int, int>;
using pll = pair<ll, ll>;
template <class T>
using oset = tree<T, null_type, less<T>,
                  rb_tree_tag, tree_order_statistics_node_update>;

ll gcd(ll a, ll b) {return a == 0? b: gcd(b%a,a);}
ll lcm(ll a, ll b) {return a * (b / gcd(a,b));}

int tc;

#define error(...) dbgh(#__VA_ARGS__, __VA_ARGS__)

template <typename T>
void dbgh(const char *names, T value)
{
  cerr << names << " = " << value << endl;
}

template <typename T, typename... Args>
void dbgh(const char *names, T value, Args... args)
{
  const char *comma = strchr(names, ',');
  cerr.write(names, comma - names) << " = " << value << endl;
  dbgh(comma + 1, args...);
}

constexpr ll INF_LL = LONG_LONG_MAX;
constexpr int INF_INT = INT_MAX;

void solve() {}

int main()
{
  ios_base::sync_with_stdio(0);
  cin.tie(0);
  int t = 1;
  // cin >> t;
  for (tc = 1; tc <= t; tc++)
    solve();
  return 0;
}
\end{minted}

	\needspace{2\baselineskip}
	{
		\phantomsection
		\subsection*{templatePy}
		\addcontentsline{toc}{subsection}{templatePy}
	}
	\begin{minted}{py}
import sys

def main():
    data = sys.stdin.read().strip().split()
    numbers = list(map(int, data))
    
    output = []
    for num in numbers:
        output.append(f"Número leído: {num}")
    
    sys.stdout.write("\n".join(output))

if __name__ == '__main__':
    main()

\end{minted}

	\needspace{5\baselineskip}
	{
		\phantomsection
		\section*{Algebra}
		\markboth{ALGEBRA}{}
		\addcontentsline{toc}{section}{Algebra}
	}
	{
		\phantomsection
		\subsection*{Berlekamp Massey}
		\addcontentsline{toc}{subsection}{Berlekamp Massey}
	}
	\begin{minted}{cpp}
/**
 * Description: Recovers a nth-order linear recurrence relation from the first
 *   2n terms. Solve can be sped up to O(n log n log k) with FFT as mul.
 * Source: https://codeforces.com/blog/entry/61306,
 *   https://github.com/kth-competitive-programming/kactl
 * Verification: https://codeforces.com/contest/506/problem/E
 * Time: O(n^2) BM, O(n^2 log k) solve
 */

template<typename T>
vector<T> berlekampMassey(const vector<T> &s) {
    int n = (int) s.size(), l = 0, m = 1;
    vector<T> b(n), c(n);
    T ld = b[0] = c[0] = 1;
    for (int i=0; i<n; i++, m++) {
        T d = s[i];
        for (int j=1; j<=l; j++)
            d += c[j] * s[i-j];
        if (d == 0)
            continue;
        vector<T> temp = c;
        T coef = d / ld;
        for (int j=m; j<n; j++)
            c[j] -= coef * b[j-m];
        if (2 * l <= i) {
            l = i + 1 - l;
            b = temp;
            ld = d;
            m = 0;
        }
    }
    c.resize(l + 1);
    c.erase(c.begin());
    for (T &x : c)
        x = -x;
    return c;
}

template<typename T>
T solve(const vector<T> &c, const vector<T> &s, long long k) {
    int n = (int) c.size();
    assert(c.size() <= s.size());

    auto mul = [&] (const vector<T> &a, const vector<T> &b) -> vector<T> {
        vector<T> ret(a.size() + b.size() - 1);
        for (int i=0; i<(int)a.size(); i++)
            for (int j=0; j<(int)b.size(); j++)
                ret[i+j] += a[i] * b[j];
        for (int i=(int)ret.size()-1; i>=n; i--)
            for (int j=n-1; j>=0; j--)
                ret[i-j-1] += ret[i] * c[j];
        ret.resize(min((int) ret.size(), n));
        return ret;
    };

    vector<T> a = n == 1 ? vector<T>{c[0]} : vector<T>{0, 1}, x{1};
    for (; k>0; k/=2) {
        if (k % 2)
            x = mul(x, a);
        a = mul(a, a);
    }
    x.resize(n);

    T ret = 0;
    for (int i=0; i<n; i++)
        ret += x[i] * s[i];
    return ret;
}

\end{minted}

	\needspace{2\baselineskip}
	{
		\phantomsection
		\subsection*{mod inv}
		\addcontentsline{toc}{subsection}{mod inv}
	}
	\begin{minted}{cpp}
#include <iostream>
using namespace std;

int euclides_extendido(int a, int b, int &x, int &y) {
    if (b == 0) {
        x = 1;
        y = 0;
        return a;
    }
    int x1, y1;
    int d = euclides_extendido(b, a % b, x1, y1);
    x = y1;
    y = x1 - (a / b) * y1;
    return d;
}

// Devuelve el inverso modular de a módulo m, o -1 si no existe
int inverso_modular(int a, int m) {
    int x, y;
    int g = euclides_extendido(a, m, x, y);
    if (g != 1) return -1; // No existe inverso si a y m no son coprimos
    x = (x % m + m) % m;   // Asegura que x sea positivo
    return x;
}


\end{minted}

	\needspace{2\baselineskip}
	{
		\phantomsection
		\subsection*{recurrencias}
		\addcontentsline{toc}{subsection}{recurrencias}
	}
	\begin{minted}{cpp}
#include <bits/stdc++.h>
using namespace std;

const int MOD = 1e9 + 7;
typedef long long ll;
ll sn, sn_1;

vector<vector<ll>> multiply(const vector<vector<ll>>& A, const vector<vector<ll>>& B) {
    int N = A.size();
    vector<vector<ll>> C(N, vector<ll>(N, 0)); 
    for (int i = 0; i < N; i++) {
        for (int j = 0; j < N; j++) {
            for (int k = 0; k < N; k++) {
                ll r = A[i][k] * B[k][j] % MOD;
                C[i][j] = (C[i][j] + r) % MOD;
            }
        }
    }
    return C;
}

vector<vector<ll>> matrixExpo(vector<vector<ll>> T, ll p) {
    int N = T.size();
    vector<vector<ll>> res(N, vector<ll>(N, 0)); 
    for (int i = 0; i < N; i++) res[i][i] = 1; 
    while (p > 0) {
        if (p & 1)
            res = multiply(res, T);
        T = multiply(T, T); 
        p >>= 1;
    }
    return res;
}

void calculateFib(int sk){
    vector<vector<ll>> F = {{1, 1}, 
                            {1, 0}};
    vector<vector<ll>> base = {{sn}, 
                                {sn_1}};
    vector<vector<ll>> res = multiply(matrixExpo(F, sk), base);
    sn= res[0][0];
    sn_1 = res[1][0];
}


\end{minted}

	\needspace{5\baselineskip}
	{
		\phantomsection
		\section*{Bits}
		\markboth{BITS}{}
		\addcontentsline{toc}{section}{Bits}
	}
	{
		\phantomsection
		\subsection*{Bits operations}
		\addcontentsline{toc}{subsection}{Bits operations}
	}
	\begin{minted}{cpp}
// Ver si n es divisible por 2^k
bool isDivisibleByPowerOf2(int n, int k) {
  int powerOf2 = 1 << k;
  return (n & (powerOf2 - 1)) == 0;
}

bool isPowerOfTwo(unsigned int n) {
  return n && !(n & (n - 1));
}

int countSetBits(int n)
{
  int count = 0;
  while (n)
  {
      n = n & (n - 1);
      count++;
  }
  return count;
}

// Cuenta el número de bits encendidos hasta el número n 
int countSetBits(int n) {
  int count = 0;
  while (n > 0) {
      int x = std::bit_width(n) - 1;
      count += x << (x - 1);
      n -= 1 << x;
      count += n + 1;
  }
  return count;
}

// Encender el i-ésimo bit
n |= (1 << i);

// Apagar el i-ésimo bit
n &= ~(1 << i);

// Alternar el i-ésimo bit
n ^= (1 << i);

// Apagar el bits más a la derecha 
n = n & (n-1); 

// Limpia todos los "trailing ones"
// 0011 0111 -> 0011 0000
n = n & (n+1)

// Enciende el último bit apagado 
// 0011 0101 -> 0011 0111
n = n | (n+1) 

// Extrae el último bit prendido
// 0011 0100 -> 0000 0100 
n = n & ~n 

// Número de bits encendidos
__builtin_popcount(x); 
__builtin_popcountll(x); 

// Posición del primer 1 desde la derecha (1-indexado)
__builtin_ffs(x);     
__builtin_ffsll(x); 

// 1 si número impar de bits encendidos
__builtin_parity(x); 
__builtin_parityll(x); 

// Nota:  clzll y ctzll no son seguras si x == 0

// Ceros a la izquierda después del primer bit encendido desde la derecha
//__builtin_clz(00000000 00000000 00000000 00010000) = 27
__builtin_clz(x); 
__builtin_clzll(x); 

// Ceros a la derecha despu+es del primer bit encendido desde la derecha 
// __builtin_ctz(00000000 00000000 00000000 00010000) = 4 
__builtin_ctz(x); 
__builtin_ctzll(x); 

\end{minted}

	\needspace{5\baselineskip}
	{
		\phantomsection
		\section*{Number Theory}
		\markboth{NUMBER THEORY}{}
		\addcontentsline{toc}{section}{Number Theory}
	}
	{
		\phantomsection
		\subsection*{Extended gcd}
		\addcontentsline{toc}{subsection}{Extended gcd}
	}
	\begin{minted}{cpp}
ll extended_gcd(ll a, ll b, ll& x, ll& y){
  if (b == 0){
    x = 1; y = 0;
    return a;
  }
  ll x1,y1;
  ll d = extended_gcd(b, a % b,x1,y1);
  x = y1;
  y = x1 - y1 * (a / b);
  return d;
}
\end{minted}

	\needspace{2\baselineskip}
	{
		\phantomsection
		\subsection*{Inverse mod}
		\addcontentsline{toc}{subsection}{Inverse mod}
	}
	\begin{minted}{cpp}
ll inverseMod(ll n, ll p){
    return binpow(n,p-2,p);
}

\end{minted}

	\needspace{2\baselineskip}
	{
		\phantomsection
		\subsection*{Is prime for large numbers}
		\addcontentsline{toc}{subsection}{Is prime for large numbers}
	}
	\begin{minted}{cpp}
// Test de primalidad de Fermat (probabilistico)
bool isPrime(ll n, int iter = 5)
{
  if (n < 4) return n == 2 || n == 3;
  if (n % 2 == 0 || n % 3 == 0) return false;
  for (int i = 0; i < iter; i++)
  {
    ll a = 2 + rand() % (n - 3);
    if (modpow(a, n - 1, n) != 1) return false;
  }
  return true;
}
\end{minted}

	\needspace{2\baselineskip}
	{
		\phantomsection
		\subsection*{Modulo long long}
		\addcontentsline{toc}{subsection}{Modulo long long}
	}
	\begin{minted}{cpp}
#include <bits/stdc++.h>
using namespace std;
using ll = long long;
constexpr ll MODL = 998244353ll;
struct mll
{
    ll x;
    mll() : x(0) {}
    mll(ll x_) : x((x_ % MODL) + (x_<0?MODL:0)) {}
    inline mll operator+(mll y) const { return mll(x + y.x); }
    inline mll operator*(mll y) const { return mll(x * y.x); }
    inline mll& operator+=(mll y) { x = (x + y.x) % MODL; return *this; }
    inline mll& operator*=(mll y) { x = (x * y.x) % MODL; return *this; }
    inline mll operator/(mll y) const { return mll(x * (y^(MODL-2)).x); }
    inline mll operator^(ll y) const
    {
        mll r = 1;
        mll a = mll(x);
        while (y > 0)
        {
            if (y & 1) r = r * a;
            a *= a;
            y >>= 1;
        }
        return r;
    }
};
ostream& operator<<(ostream& os, mll x) { os << x.x; return os; }
istream& operator>>(istream& is, mll& x) { ll y; is >> y; x = y; return is; }

\end{minted}

	\needspace{2\baselineskip}
	{
		\phantomsection
		\subsection*{NTT}
		\addcontentsline{toc}{subsection}{NTT}
	}
	\begin{minted}{cpp}
// Transformada discreta de Fourier en Z_p
const int mod = 7340033;
const int root = 5; // (2^k)-ésima raíz de la unidad
const int root_1 = 4404020; // inversa de root (modulo mod)
const int root_pw = 1 << 20; // 2^k

// para p primo de la forma p=c2^k + 1, g^c es (2^k)-ésima raíz de la unidad,
// donde g es raíz primitiva de p, i.e. todo factor p_i de (p-1) cumple que
// g^((p-1)/p_i) mod p es congruente a 1.

// FFT + inverse FFT sobre Z_p
void fft(vector<int> & a, bool invert) {
    int n = a.size();

    for (int i = 1, j = 0; i < n; i++) {
        int bit = n >> 1;
        for (; j & bit; bit >>= 1)
            j ^= bit;
        j ^= bit;

        if (i < j)
            swap(a[i], a[j]);
    }

    for (int len = 2; len <= n; len <<= 1) {
        int wlen = invert ? root_1 : root;
        for (int i = len; i < root_pw; i <<= 1)
            wlen = (int)(1LL * wlen * wlen % mod);

        for (int i = 0; i < n; i += len) {
            int w = 1;
            for (int j = 0; j < len / 2; j++) {
                int u = a[i+j], v = (int)(1LL * a[i+j+len/2] * w % mod);
                a[i+j] = u + v < mod ? u + v : u + v - mod;
                a[i+j+len/2] = u - v >= 0 ? u - v : u - v + mod;
                w = (int)(1LL * w * wlen % mod);
            }
        }
    }

    if (invert) {
        int n_1 = inverse(n, mod); // TODO: inverso modular de n
        for (int & x : a)
            x = (int)(1LL * x * n_1 % mod);
    }
}

\end{minted}

	\needspace{2\baselineskip}
	{
		\phantomsection
		\subsection*{Prime factorization}
		\addcontentsline{toc}{subsection}{Prime factorization}
	}
	\begin{minted}{cpp}

vector<ll> trial_division(ll n) {
  vector<ll> factorization;
  for (int d : {2, 3, 5}) {
    while (n % d == 0) {
      factorization.push_back(d);
      n /= d;
    }
  }
  static array<int, 8> increments = {4, 2, 4, 2, 4, 6, 2, 6};
  int i = 0;
  for (ll d = 7; d * d <= n; d += increments[i++]) {
    while (n % d == 0) {
      factorization.push_back(d);
      n /= d;
    }
    if (i == 8)
      i = 0;
  }
  if (n > 1)
    factorization.push_back(n);

  // Si factorization.size() == 1 entonces n es primo
  return factorization;
}
\end{minted}

	\needspace{5\baselineskip}
	{
		\phantomsection
		\section*{Data Structures}
		\markboth{DATA STRUCTURES}{}
		\addcontentsline{toc}{section}{Data Structures}
	}
	{
		\phantomsection
		\subsection*{AVL}
		\addcontentsline{toc}{subsection}{AVL}
	}
	\begin{minted}{cpp}
using ll = long long;
using i16 = short int;

// nodo avl
struct avl_n
{
  avl_n *l, *r; // izq/der
  i16 h; // altura
  ll sz; // tamaño subárbol
  int x; // valor

  // constructor
  avl_n(int x) : l(0), r(0), h(1), sz(1), x(x)
  { }

  // computa balance
  inline i16 balance() { return (r ? r->h : 0) - (l ? l->h : 0); }

  // encuentra el mínimo del subárbol
  avl_n *min()
  {
    avl_n *r = this;
    while (r->l)
      r = r->l;
    return r;
  }

  // encuentra el máximo del subárbol
  avl_n *max()
  {
    avl_n *r = this;
    while (r->r)
      r = r->r;
    return r;
  }
};

// arbol avl
struct avl_t
{
  // raíz
  avl_n *r;

  // constructor
  avl_t() : r(0)
  { }

  // auxiliar para computar el tamaño de un nodo posiblemente nulo
  inline int v_sz(avl_n *&v) { return v ? v->sz : 0; }
  // auxiliar para computar la altura de un nodo posiblemente nulo
  inline int v_h(avl_n *&v) { return v ? v->h : 0; }

  // actualiza altura y tamaño de un subárbol (vértice)
  inline void update(avl_n *&v)
  {
    if (!v)
      return;
    v->h = 1 + max(v_h(v->l), v_h(v->r));
    v->sz = 1 + v_sz(v->l) + v_sz(v->r);
  }

  // obtiene el tamaño del árbol
  inline int sz() { return v_sz(r); }

  // giro izq
  void lrot(avl_n *&v)
  {
    avl_n *y = v->r, *t = y->l;
    y->l = v, v->r = t;
    update(v), update(y);
    v = y;
  }

  // giro der
  void rrot(avl_n *&v)
  {
    avl_n *x = v->l, *t = x->r;
    x->r = v, v->l = t;
    update(v), update(x);
    v = x;
  }

  // actualiza el balance y realiza las rotaciones para mantenerlo
  void updateb(avl_n *&v)
  {
    if (!v)
      return;
    i16 b = v->balance();
    if (b > 1)
    {
      if (v->r->balance() < 0)
        rrot(v->r);
      lrot(v);
    }
    else if (b < -1)
    {
      if (v->l->balance() > 0)
        lrot(v->l);
      rrot(v);
    }
  }

  // inserta un nuevo valor al arbol
  void insert(avl_n *&v, int x)
  {
    if (v)
    {
      if (x < v->x)
        insert(v->l, x);
      else
        insert(v->r, x);
      update(v), updateb(v);
    }
    else
    {
      v = new avl_n(x);
    }
  }

  // busca si el valor se encuentra en el árbol
  avl_n *search(int &x)
  {
    avl_n *v = r;
    while (v)
    {
      if (x == v->x)
        break;
      v = (x < v->x ? v->l : v->r);
    }
    return v;
  }

  // borra el valor del árbol si existe
  void erase(avl_n *&v, int &x)
  {
    if (!v)
      return;
    if (x < v->x)
      erase(v->l, x);
    else if (x > v->x)
      erase(v->r, x);
    else
    {
      if (!v->l)
        v = v->r;
      else if (!v->r)
        v = v->l;
      else
      {
        v->x = v->r->min()->x;
        erase(v->r, v->x);
      }
    }
    update(v), updateb(v);
  }

  // inserta un valor al árbol
  void insert(int x)
  {
    insert(r, x);
  }

  // elimina un valor del árbol
  void erase(int x)
  {
    erase(r, x);
  }

  // encuentra el k-ésimo elemento
  avl_n *kth(int i)
  {
    avl_n *v = r;
    while (i != v_sz(v->l))
    {
      if (i < v_sz(v->l))
      {
        v = v->l;
      }
      else
      {
        i -= v_sz(v->l) + 1;
        v = v->r;
      }
    }
    return v;
  }

  // imprime la mediana (o su promedio si hay dos, o indica que no hay elementos
  // suficientes para una mediana)
  void print_median()
  {
    if (sz() == 0)
    {
      cout << "Empty!\n";
    }
    else if (sz() % 2)
    {
      cout << kth(sz() / 2)->x << '\n';
    }
    else
    {
      avl_n *upmed = kth(sz() >> 1);
      avl_n *downmed = (upmed->l ? upmed->l->max() : kth((sz() >> 1) - 1));
      cout << (downmed->x + ((upmed->x - downmed->x) >> 1)) << '\n';
    }
  }
};

\end{minted}

	\needspace{2\baselineskip}
	{
		\phantomsection
		\subsection*{Bitset}
		\addcontentsline{toc}{subsection}{Bitset}
	}
	\begin{minted}{cpp}
// Tamaño del bitset (usa hasta 32 si quieres convertir a int sin problemas)
const int N = 8;
int main() {
    // 🔹 Crear bitsets
    bitset<N> bs1(string("10101001"));   // Desde string binario
    int num = 42;
    bitset<N> bs2(num);                  // Desde entero (42 = 00101010)

    cout << "bs1: " << bs1 << endl;
    cout << "bs2: " << bs2 << endl;

    // 🔹 Iterar sobre bits
    cout << "Bits activos (1) en bs1: ";
    for (int i = 0; i < N; ++i) {
        if (bs1[i]) cout << i << " ";
    }
    cout << "\n";

    // 🔹 Operaciones lógicas entre bitsets
    cout << "bs1 & bs2: " << (bs1 & bs2) << endl;
    cout << "bs1 | bs2: " << (bs1 | bs2) << endl;
    cout << "bs1 ^ bs2: " << (bs1 ^ bs2) << endl;

    // 🔹 Métodos importantes
    cout << "bs1.any():  " << bs1.any() << "  // ¿Algún bit en 1?\n";
    cout << "bs1.none(): " << bs1.none() << " // ¿Todos los bits en 0?\n";
    cout << "bs1.all():  " << bs1.all() << "  // ¿Todos los bits en 1?\n";
    cout << "bs1.count(): " << bs1.count() << " // ¿Cuántos bits en 1?\n";

    // 🔹 Modificación de bits
    bs1.set(2);     // Pone a 1 el bit 2
    bs1.reset(0);   // Pone a 0 el bit 0
    bs1.flip(1);    // Invierte el bit 1

    cout << "bs1 modificado: " << bs1 << endl;

    bs1.set();      // Todos los bits a 1
    cout << "bs1.set(): " << bs1 << endl;

    bs1.reset();    // Todos los bits a 0
    cout << "bs1.reset(): " << bs1 << endl;

    bs1.flip();     // Invierte todos los bits
    cout << "bs1.flip(): " << bs1 << endl;

    // 🔹 Conversión a entero
    unsigned long val_ul = bs2.to_ulong();         // Hasta 32 bits
    unsigned long long val_ull = bs2.to_ullong();  // Hasta 64 bits
    int val_int = static_cast<int>(val_ul);        // Seguro si N <= 32

    cout << "bs2 como unsigned long: " << val_ul << endl;
    cout << "bs2 como int: " << val_int << endl;

    // 🔹 Convertir int a bitset nuevamente
    bitset<N> fromInt(val_int);
    cout << "Convertido desde int: " << fromInt << endl;

    return 0;
}

\end{minted}

	\needspace{2\baselineskip}
	{
		\phantomsection
		\subsection*{Deque}
		\addcontentsline{toc}{subsection}{Deque}
	}
	\begin{minted}{cpp}
int main() {
    deque<int> dq;

    // Insertar elementos al final de la deque (push_back)
    dq.push_back(10);
    dq.push_back(20);
    dq.push_back(30);
    dq.push_back(40);

    // Insertar elementos al frente de la deque (push_front)
    dq.push_front(5);
    dq.push_front(2);

    // Mostrar los elementos de la deque
    cout << "Elementos en la deque:" << endl;
    for (int el : dq) {
        cout << el << " ";
    }
    cout << endl;

    // Eliminar el elemento al frente (pop_front)
    dq.pop_front();

    // Eliminar el elemento al final (pop_back)
    dq.pop_back();

    cout << "Elementos después de eliminar el frente y el final:" << endl;
    for (int el : dq) {
        cout << el << " ";
    }
    cout << endl;

    // Verificar si la deque está vacía
    if (dq.empty()) {
        cout << "La deque está vacía." << endl;
    } else {
        cout << "La deque no está vacía." << endl;
    }

    // Mostrar el primer y último elemento
    cout << "El primer elemento es: " << dq.front() << endl;
    cout << "El último elemento es: " << dq.back() << endl;
    return 0;
}

\end{minted}

	\needspace{2\baselineskip}
	{
		\phantomsection
		\subsection*{DSU}
		\addcontentsline{toc}{subsection}{DSU}
	}
	\begin{minted}{cpp}
typedef vector<int> vi;
/**
 * Estructura DSU (Disjoint Set Union) o Union-Find.
 * Permite mantener conjuntos disjuntos con operaciones eficientes.
 */
struct DSU
{
  vi parent;   // parent[i]: padre del nodo i
  vi rank;     // rank[i]: estimación de la profundidad del árbol del conjunto que contiene a i
  vi setSize;  // setSize[i]: tamaño del conjunto cuya raíz es i
  int numSets; // Número actual de conjuntos disjuntos

  // Constructor: Inicializa DSU con N elementos (1 a N)
  DSU(int N)
  {
    N++; // Para trabajar con índices desde 1
    parent.resize(N);
    rank.assign(N, 0);
    setSize.assign(N, 1);
    for (int i = 1; i < N; ++i)
      parent[i] = i;
    numSets = N - 1;
  }

  // Devuelve la raíz del conjunto que contiene a 'i' (con compresión de caminos)
  int find(int i)
  {
    if (parent[i] != i)
      parent[i] = find(parent[i]);
    return parent[i];
  }

  // Retorna true si 'i' y 'j' están en el mismo conjunto
  bool isSameSet(int i, int j)
  {
    return find(i) == find(j);
  }

  // Retorna el número de conjuntos disjuntos actuales
  int countSets()
  { return numSets; }

  // Retorna el tamaño del conjunto que contiene al elemento 'i'
  int size(int i)
  { return setSize[find(i)]; }

  // Une los conjuntos que contienen a 'i' y 'j'
  void unify(int i, int j)
  {
    if (isSameSet(i, j))
      return;
    int x = find(i), y = find(j);

    // Unión por rango para mantener árboles bajos
    if (rank[x] > rank[y])
      swap(x, y);

    parent[x] = y;
    if (rank[x] == rank[y])
      rank[y]++;

    setSize[y] += setSize[x];
    --numSets;
  }
};

int main()
{
  // Ejemplo de uso
  DSU dsu(5); // Elementos del 1 al 5

  dsu.unify(1, 2);
  dsu.unify(3, 4);

  cout << dsu.isSameSet(1, 2) << endl; //  1 (true)
  cout << dsu.isSameSet(1, 3) << endl; //  0 (false)
  cout << dsu.size(1) << endl;         //  2
  cout << dsu.countSets() << endl;     //  3

  dsu.unify(2, 3);
  cout << dsu.countSets() << endl; //  2

  // Ejemplo componentes conexas
  int n = 6; // nodos del 1 al 6
  vector<pair<int, int>> edges = {
      {1, 2}, {2, 3}, {4, 5}};

  DSU dsu(n);

  // Unir nodos conectados
  for (auto &[u, v] : edges)
    dsu.unify(u, v);

  // Mostrar cuántas componentes conexas hay
  cout << "Componentes conexas: " << dsu.countSets() << endl; // Debería imprimir 3

  // Mostrar tamaño de la componente de cada nodo
  for (int i = 1; i <= n; ++i)
  {
    cout << "Nodo " << i << " pertenece a componente de tamaño: " << dsu.size(i) << endl;
  }
}

\end{minted}

	\needspace{2\baselineskip}
	{
		\phantomsection
		\subsection*{List}
		\addcontentsline{toc}{subsection}{List}
	}
	\begin{minted}{cpp}
#include <iostream>
#include <list>
#include <iterator> // Necesario para next() y prev()
using namespace std;

int main() {
    list<int> lst;
    lst.push_back(10);
    lst.push_back(20);
    lst.push_back(30);
    lst.push_back(40);
    lst.push_back(50);

    cout << "Elementos en la lista (usando iteradores):" << endl;
    for (auto it = lst.begin(); it != lst.end(); ++it) {
        cout << *it << " ";
    }
    cout << endl;

    auto it = lst.begin();
    cout << "El primer elemento es: " << *it << endl;
    it = next(it, 2); // Avanzar el iterador 2 posiciones
    cout << "El tercer elemento (después de avanzar 2 pasos) es: " << *it << endl;

    // Uso de `prev()` para retroceder el iterador
    it = prev(it, 1); // Retroceder 1 posición
    cout << "El segundo elemento (después de retroceder 1 paso) es: " << *it << endl;

    // Inserción de un elemento antes de un iterador específico
    it = next(lst.begin(), 2); // Moverse al tercer elemento
    lst.insert(it, 35); // Insertar 35 antes del tercer elemento

    // Mostrar los elementos después de la inserción
    cout << "\nElementos después de insertar 35 antes del tercer elemento:" << endl;
    for (auto it = lst.begin(); it != lst.end(); ++it) {
        cout << *it << " ";
    }
    cout << endl;

    // Eliminar un elemento usando un iterador
    it = next(lst.begin(), 3); // Moverse al cuarto elemento
    lst.erase(it); // Eliminar el cuarto elemento (40)

    // Acceder al primer y último elemento usando `front()` y `back()`
    cout << "\nEl primer elemento es: " << lst.front() << endl;
    cout << "El último elemento es: " << lst.back() << endl;
    return 0;
}

\end{minted}

	\needspace{2\baselineskip}
	{
		\phantomsection
		\subsection*{MinMax Heap}
		\addcontentsline{toc}{subsection}{MinMax Heap}
	}
	\begin{minted}{cpp}
// Ojo. las operaciones de lectura/eliminacion asumen que si hay cosas
#include <bits/stdc++.h>
using namespace std;
struct MMH{
    vector<int> h;
    bool minLv(int i){int d=0;for(;i;i=(i-1)/2)d++;return d%2==0;}
    void push(int x){
        h.push_back(x);int i=h.size()-1,p=(i-1)/2;
        if((minLv(i)&&h[i]>h[p])||(!minLv(i)&&h[i]<h[p]))
            swap(h[i],h[p]),minLv(i)?upMax(p):upMin(p);
        else minLv(i)?upMin(i):upMax(i);
    }
    int getMin(){return h[0];}
    int getMax(){return h.size()<2?h[0]:(h.size()==2?h[1]:max(h[1],h[2]));}
    void popMin(){h[0]=h.back();h.pop_back();down(0);}
    void popMax(){int m=h.size()==2||h[1]>h[2]?1:2;
        h[m]=h.back();h.pop_back();down(m);}
    void upMin(int i){while(i>=3){int g=((i-1)/2-1)/2;
        if(h[i]<h[g])swap(h[i],h[g]),i=g;else break;}}
    void upMax(int i){while(i>=3){int g=((i-1)/2-1)/2;
        if(h[i]>h[g])swap(h[i],h[g]),i=g;else break;}}
    void down(int i){minLv(i)?downMin(i):downMax(i);}
    void downMin(int i){for(;;){int m=-1;
        for(int d=1;d<=4;d++){int c=2*i+d;
            if(c<h.size()&&(m==-1||h[c]<h[m]))m=c;}
        if(m==-1)break;
        if(m>=4*i+3){if(h[m]<h[i]){swap(h[m],h[i]);
            int p=(m-1)/2;if(h[m]>h[p])swap(h[m],h[p]);downMin(m);}}
        else{if(h[m]<h[i])swap(h[m],h[i]);else break;}}}
    void downMax(int i){for(;;){int m=-1;
        for(int d=1;d<=4;d++){int c=2*i+d;
            if(c<h.size()&&(m==-1||h[c]>h[m]))m=c;}
        if(m==-1)break;
        if(m>=4*i+3){if(h[m]>h[i]){swap(h[m],h[i]);
            int p=(m-1)/2;if(h[m]<h[p])swap(h[m],h[p]);downMax(m);}}
        else{if(h[m]>h[i])swap(h[m],h[i]);else break;}}}
};


\end{minted}

	\needspace{2\baselineskip}
	{
		\phantomsection
		\subsection*{Monotonic MeteSaca}
		\addcontentsline{toc}{subsection}{Monotonic MeteSaca}
	}
	\begin{minted}{cpp}
struct mono_stack
{
    vector<pair<int,int>> st;

    void push(int v, int idx)
    {
        // cambiar <= a < si quieres ser estricto
        while (!st.empty() && st.back().first <= v)
            st.pop_back();
        st.emplace_back(v, idx);
    }

    void pop()
    {
        if (!st.empty())
            st.pop_back();
    }

    int top() const // verificar primero empty()!!!
    {
        return st.back().first;
    }

    bool empty() const { return st.empty(); }
    size_t size() const { return st.size(); }
};

struct mono_queue
{
    deque<pair<int,int>> dq; // {valor, idx}

    void push(int v, int idx)
    {
        // cambiar <= a < si quieres ser estricto
        while (!dq.empty() && dq.back().first <= v)
            dq.pop_back();
        dq.push_back({v,idx});
    }

    void pop(int idx_out)
    {
        // cambiar <= a < si quieres ser estricto
        if (!dq.empty() && dq.front().second <= idx_out)
            dq.pop_front();
    }

    int get() const // verificar primero empty()!!!
    {
        return dq.front().first;
    }
};

// Comparadores estándar: less<T>, less_equal<T>, greater<T>, greater_equal<T>
template<typename T, typename Compare = less<T>>
struct mono_deque
{
    deque<T> dq;
    Compare comp;

    void push_back_mono(const T& val)
    {
        while (!dq.empty() && comp(dq.back(), val))
            dq.pop_back();
        dq.push_back(val);
    }

    void push_front_mono(const T& val)
    {
        while (!dq.empty() && comp(val, dq.front()))
            dq.pop_front();
        dq.push_front(val);
    }

    void pop_back() {dq.pop_back();} // verificar empty()!!!
    void pop_front() {dq.pop_front();} // verificar empty()!!!
    const T& front() const {return dq.front();} // verificar empty()!!!
    const T& back() const {return dq.back();} // verificar empty()!!!
    bool empty() const {return dq.empty();}
    size_t size() const {return dq.size();}
};

\end{minted}

	\needspace{2\baselineskip}
	{
		\phantomsection
		\subsection*{Ordered Set}
		\addcontentsline{toc}{subsection}{Ordered Set}
	}
	\begin{minted}{cpp}
#include <bits/stdc++.h>
#include <ext/pb_ds/assoc_container.hpp>
#include <ext/pb_ds/tree_policy.hpp>
using namespace std;
using namespace __gnu_pbds;

// Se puede usar como un set, tiene todas sus funcionalidades pero agrega dos funciones extra
template <class T>
using oset = tree<T, null_type, less<T>,
                  rb_tree_tag, tree_order_statistics_node_update>;
// find_by_order, order_of_key
// También en vez de less<T> se puede poner greater<T> mayor estricto (>)

int main()
{
  oset<int> A; // declaración

  // Insertar elementos
  A.insert(1);
  A.insert(10);
  A.insert(2);
  A.insert(7);
  A.insert(2); // Solo contiene valores unicos

  // A contains = 1 2 7 10
  cout << "A = ";
  for (auto i : A)
    cout << i << " ";
  cout << endl;

  // Encontrar el kth elemento
  cout << "0th element: " << *A.find_by_order(0) << endl; //  1
  cout << "1st element: " << *A.find_by_order(1) << endl; //  2
  cout << "2nd element: " << *A.find_by_order(2) << endl; //  7
  cout << "3rd element: " << *A.find_by_order(3) << endl; //  10
  auto invalido = A.find_by_order(10);
  if (invalido == A.end())
    cout << "No existe ese elemento" << endl;
  cout << endl;

  // Encontrar el número de elementos menores que x
  cout << "No. of elems smaller than 6: " << A.order_of_key(6) << endl;   //  2
  cout << "No. of elems smaller than 11: " << A.order_of_key(11) << endl; //  4
  cout << "No. of elems smaller than 1: " << A.order_of_key(1) << endl;   //  0
  cout << endl;

  // lower bound -> Primer elemento >= que x
  cout << "Lower Bound of 6: " << *A.lower_bound(6) << endl;
  cout << "Lower Bound of 2: " << *A.lower_bound(2) << endl;
  cout << endl;

  // Upper bound -> Primer elemento > que x
  cout << "Upper Bound of 6: " << *A.upper_bound(6) << endl;
  cout << "Upper Bound of 1: " << *A.upper_bound(1) << endl;
  cout << endl;

  // Eliminar elementos
  A.erase(1);
  A.erase(11); // Eliminar un elemento que no está no hace nada

  // A contains = 2 7 10
  cout << "A = ";
  for (auto i : A)
    cout << i << " ";
  cout << endl;
}
\end{minted}

	\needspace{2\baselineskip}
	{
		\phantomsection
		\subsection*{Priority Queue}
		\addcontentsline{toc}{subsection}{Priority Queue}
	}
	\begin{minted}{cpp}
struct Comparador
{
  bool operator()(const pair<int, string> &a, const pair<int, string> &b)
  {
    return a.first > b.first; // Ordenar por el primer valor del par (min-heap por valor de entero)
  }
};

int main()
{
  // Crear una priority_queue con el comportamiento por defecto (max-heap)
  priority_queue<int> pqMax;

  // Insertar elementos en el max-heap
  pqMax.push(5);
  pqMax.push(2);
  pqMax.push(8);
  pqMax.push(1);
  pqMax.push(3);

  cout << "Elementos en el max-heap (prioridad alta a baja):" << endl;

  // Extraer elementos (se extrae el de mayor valor primero)
  // 8 5 3 2 1
  while (!pqMax.empty())
  {
    cout << pqMax.top() << " "; // Muestra el elemento de mayor prioridad (en max-heap, el más grande)
    pqMax.pop();                // Elimina el elemento de mayor prioridad
  }
  cout << endl;

  // Crear una priority_queue con el comportamiento de min-heap
  priority_queue<int, vector<int>, greater<int>> pqMin;

  // Insertar elementos en la min-heap
  pqMin.push(5);
  pqMin.push(2);
  pqMin.push(8);
  pqMin.push(1);
  pqMin.push(3);

  cout << "Elementos en el min-heap (prioridad baja a alta):" << endl;

  // Extraer elementos (se extrae el de menor valor primero)
  // 1 2 3 5 8
  while (!pqMin.empty())
  {
    cout << pqMin.top() << " "; // Muestra el elemento de mayor prioridad (en max-heap, el más grande)
    pqMin.pop();                // Elimina el elemento de mayor prioridad
  }
  cout << endl;

  // Priority Queue con comparador persoanlizado
  priority_queue<pair<int, string>, vector<pair<int, string>>, Comparador> pq;

  pq.push({5, "cinco"});
  pq.push({2, "dos"});
  pq.push({8, "ocho"});
  pq.push({1, "uno"});
  pq.push({3, "tres"});

  while (!pq.empty())
  {
    cout << pq.top().first << " - " << pq.top().second << endl;
    pq.pop();
  }
  // 1 - uno
  // 2 - dos ...
  return 0;
}

\end{minted}

	\needspace{3\baselineskip}
	{
		\phantomsection
		\subsection*{LazySegmentTree}
		\addcontentsline{toc}{subsection}{LazySegmentTree}
	}
	{
		\phantomsection
		\subsubsection*{Assing On A Range And Look On Position}
		\addcontentsline{toc}{subsubsection}{Assing On A Range And Look On Position}
	}
	\begin{minted}{cpp}
// Lazy segment tree para queries de asignamiento en rango
// y buscar en una posicion
template <typename T>
struct LazySegmentTree
{
  int n;
  vector<T> st, lazy;
  vector<bool> marked;

  void push(int v)
  {
    if (marked[v])
    {
      lazy[2 * v] = lazy[2 * v + 1] = lazy[v];
      marked[2 * v] = marked[2 * v + 1] = true;
      lazy[v] = 0;
      marked[v] = false;
    }
  }

  void build(vector<T> &nums, int v, int tl, int tr)
  {
    if (tl == tr) 
      st[v] = nums[tl];
    else
    {
      int mid = (tl + tr) / 2;
      build(nums, 2 * v, tl, mid);
      build(nums, 2 * v + 1, mid + 1, tr);
    }
  }

  LazySegmentTree(vector<T> &nums)
  {
    n = nums.size();
    st.resize(4 * n);
    lazy.resize(4 * n);
    marked.resize(4 * n);
    build(nums, 1, 0, n - 1);
  }

  T query(int v, int tl, int tr, int pos)
  {
    if (tl == tr) 
      return marked[v] ? lazy[v] : st[v];
    push(v);
    int mid = (tl + tr) / 2;
    if (pos <= mid)
      return query(2 * v, tl, mid, pos);
    return query(2 * v + 1, mid + 1, tr, pos);
  }

  void update(int v, int tl, int tr, int l, int r, T val)
  {
    if (tl > r || tr < l)
      return;
    if (tl >= l && tr <= r)
      lazy[v] = val, marked[v] = true;
    else
    {
      push(v);
      int mid = (tl + tr) / 2;
      update(2 * v, tl, mid, l, r, val);
      update(2 * v + 1, mid + 1, tr, l, r, val);
    }
  }
};

int main()
{
  int n, q; cin >> n >> q;
  vll nums(n);
  LazySegmentTree<ll> st(nums);
  while (q--)
  {
    int op;
    cin >> op;
    if (op == 1)
    {
      int l, r; ll v;
      cin >> l >> r >> v;
      l--; r--; // Solo restamos si las posiciones son 1-indexed
      st.update(1, 0, n-1, l, r, v);
    }
    else
    {
      int pos; cin >> pos;
      pos--; // Solo restamos si las posiciones son 1-indexed
      cout << st.query(1, 0, n-1, pos) << "\n";
    }
  }
  return 0;
}
\end{minted}

	\needspace{1\baselineskip}
	{
		\phantomsection
		\subsubsection*{Sum And Update On Range}
		\addcontentsline{toc}{subsubsection}{Sum And Update On Range}
	}
	\begin{minted}{cpp}
// Lazy Segment Tree para suma de un valor x en rangos 
// y obtener la suma de un rango [l, r]
template <typename T>
struct LazySegmentTree
{
  T neuter;
  vector<ll> st;
  vector<ll> lazy;

  LazySegmentTree(vector<T> &nums)
  {
    neuter = 0;
    int n = nums.size();
    st.assign(4 * n, neuter);
    lazy.assign(4 * n, 0);
    build(1, 0, n - 1, nums);
  }

  inline T merge(T a, T b) { return a + b; }

  void push(int v, int tl, int tr)
  {
    lazy[v * 2] += lazy[v];
    lazy[v * 2 + 1] += lazy[v];
    int tm = tl + (tr - tl) / 2;
    st[v * 2] += lazy[v] * ((tm - tl) + 1);
    st[v * 2 + 1] += lazy[v] * (tr - tm);
    lazy[v] = 0;
  }
  void build(int v, int tl, int tr, vector<T> &nums)
  {
    if (tl == tr)
    {
      st[v] = nums[tl];
      return;
    }
    int tm = tl + (tr - tl) / 2;
    build(v * 2, tl, tm, nums);
    build(v * 2 + 1, tm + 1, tr, nums);
    st[v] = merge(st[v * 2], st[v * 2 + 1]);
  }

  void update(int v, int tl, int tr, int l, int r, T val)
  {
    if (tl > r || tr < l) return;
    
    if (tl >= l && tr <= r)
    {
      st[v] += val * ((tr - tl) + 1);
      lazy[v] += val;
      return;
    }
    push(v, tl, tr);
    int tm = tl + (tr - tl) / 2;
    update(v * 2, tl, tm, l, r, val);
    update(v * 2 + 1, tm + 1, tr, l, r, val);
    st[v] = merge(st[v * 2], st[v * 2 + 1]);
  }

  T query(int v, int tl, int tr, int l, int r)
  {
    if (tl > r || tr < l) return neuter;
    

    if (tl >= l && tr <= r) return st[v];
    
    push(v, tl, tr);
    int tm = tl + (tr - tl) / 2;
    return merge(query(v * 2, tl, tm, l, r), 
           query(v * 2 + 1, tm + 1, tr, l, r));
  }
};
\end{minted}

	\needspace{1\baselineskip}
	{
		\phantomsection
		\subsubsection*{Sum On A Range And Minimum}
		\addcontentsline{toc}{subsubsection}{Sum On A Range And Minimum}
	}
	\begin{minted}{cpp}
// Lazy segment tree para querys de sumar en un rango 
// y encontrar el minimo de un rango [l, r]
template <typename T>
struct LazySegmentTree
{
  int n;
  int inf = INT_MAX;
  vector<T> st, lazy;

  void push(int v)
  {
    st[2 * v] += lazy[v];
    lazy[2 * v] += lazy[v];
    st[2 * v + 1] += lazy[v];
    lazy[2 * v + 1] += lazy[v];
    lazy[v] = 0;
  }

  void build(vector<T> &nums, int v, int tl, int tr)
  {
    if (tl == tr)
      st[v] = nums[tl];
    else
    {
      int mid = (tl + tr) / 2;
      build(nums, 2 * v, tl, mid);
      build(nums, 2 * v + 1, mid + 1, tr);
    }
  }

  LazySegmentTree(vector<T> &nums)
  {
    n = nums.size();
    st.resize(4 * n);
    lazy.resize(4 * n);
    build(nums, 1, 0, n - 1);
  }

  T query(int v, int tl, int tr, int l, int r)
  {
    if (tl > r || tr < l)
      return inf;
    if (tl >= l && tr <= r)
      return st[v];
    push(v);
    int mid = (tl + tr) / 2;
    return min(query(2 * v, tl, mid, l, r),
               query(2 * v + 1, mid + 1, tr, l, r));
  }

  void update(int v, int tl, int tr, int l, int r, T val)
  {
    if (tl > r || tr < l)
      return;
    if (tl >= l && tr <= r)
      lazy[v] += val, st[v] += val;
    else
    {
      push(v);
      int mid = (tl + tr) / 2;
      update(2 * v, tl, mid, l, r, val);
      update(2 * v + 1, mid + 1, tr, l, r, val);
    }
    st[v] = min(st[2 * v], st[2 * v + 1]);
  }
};

\end{minted}

	\needspace{3\baselineskip}
	{
		\phantomsection
		\subsection*{MergeSortTree}
		\addcontentsline{toc}{subsection}{MergeSortTree}
	}
	{
		\phantomsection
		\subsubsection*{MergeSortTree}
		\addcontentsline{toc}{subsubsection}{MergeSortTree}
	}
	\begin{minted}{cpp}
template <typename T>
struct MergeSortTree
{
  vector<vector<T>> st;
  MergeSortTree(vector<T> &nums)
  {
    st.resize(nums.size() * 4);
    build(1, 0, nums.size() - 1, nums);
  }

  void build(int v, int tl, int tr, vector<T> &nums)
  {
    if (tl == tr)
    {
      st[v] = {nums[tl]};
      return;
    }

    int tm = tl + (tr - tl) / 2;
    build(v * 2, tl, tm, nums);
    build(v * 2 + 1, tm + 1, tr, nums);

    merge(all(st[v * 2]), all(st[v * 2 + 1]), back_inserter(st[v]));
  }

  // Número de elementos iguales a val   
  T query(int v, int tl, int tr, int l, int r, T val)
  {
    if (tl > r || tr < l)
    {
      return 0;
    }

    if (tl >= l && tr <= r)
    {
      auto it = lower_bound(all(st[v]), val);
      auto jt = upper_bound(all(st[v]), val);

      if (it == st[v].end())
      {
        return 0;
      }

      int j = jt - st[v].begin();
      int i = it - st[v].begin();
      return j - i;
    }

    int tm = tl + (tr - tl) / 2;

    return query(v * 2, tl, tm, l, r, val) + query(v * 2 + 1, tm + 1, tr, l, r, val);
  }
};
\end{minted}

	\needspace{3\baselineskip}
	{
		\phantomsection
		\subsection*{SegmentTree}
		\addcontentsline{toc}{subsection}{SegmentTree}
	}
	{
		\phantomsection
		\subsubsection*{2D Segment Tree}
		\addcontentsline{toc}{subsubsection}{2D Segment Tree}
	}
	\begin{minted}{cpp}
// Primero se aplica cualquier cosa sobre X y después sobre Y. Esto aplica a
// todas las funciones siguientes, por lo que se debe llamar su versión en X.

// La entrada es una matriz a[][], y el árbol se almacena sobre t[][]

// CONSTRUCCIÓN

// Esto es interno, se debe llamar build_x el cual llamará esta función
// Build the Y-axis segment tree for a fixed X-node vx
// vx: index in X-axis segment tree
// lx, rx: X-range covered by vx
// vy: index in Y-axis segment tree for this vx
// ly, ry: Y-range covered by vy
void build_y(int vx, int lx, int rx, int vy, int ly, int ry) {
	if (ly == ry)
		if (lx == rx)
			t[vx][vy] = a[lx][ly];
		else
			t[vx][vy] = t[vx*2][vy] + t[vx*2+1][vy];
	else {
		int my = (ly + ry) / 2;
		build_y (vx, lx, rx, vy*2, ly, my);
		build_y (vx, lx, rx, vy*2+1, my+1, ry);
		t[vx][vy] = t[vx][vy*2] + t[vx][vy*2+1];
	}
}
 
// Sobre el segment-tree en X se construyen segment trees en Y
// Build the X-axis segment tree
// vx: index in X-axis segment tree
// lx, rx: X-range covered by vx
void build_x(int vx, int lx, int rx) {
	if (lx != rx) {
		int mx = (lx + rx) / 2;
		build_x (vx*2, lx, mx);
		build_x (vx*2+1, mx+1, rx);
	}
	build_y (vx, lx, rx, 1, 0, m-1);
}

// EJEMPLO DE QUERY

// Esto es interno, realmente se debería llamar sum_x, pues se procesa primero
// un strip en X y después los strips en Y.
// Query the Y-axis segment tree for a fixed X-node vx
// vx: fixed X-node index
// vy: current Y-node index
// tly, try_: total Y-range covered by vy
// ly, ry: query Y-range
int sum_y(int vx, int vy, int tly, int try_, int ly, int ry) {
	if (ly > ry)
		return 0;
	if (ly == tly && try_ == ry)
		return t[vx][vy];
	int tmy = (tly + try_) / 2;
	return sum_y (vx, vy*2, tly, tmy, ly, min(ry,tmy))
		+ sum_y (vx, vy*2+1, tmy+1, try_, max(ly,tmy+1), ry);
}
 
// O(logn logm) pues se hacen todos los strips en Y sobre el strip en X
// Query the X-axis segment tree
// vx: current X-node index
// tlx, trx: total X-range covered by vx
// lx, rx: query X-range
// ly, ry: query Y-range
int sum_x(int vx, int tlx, int trx, int lx, int rx, int ly, int ry) {
	if (lx > rx)
		return 0;
	if (lx == tlx && trx == rx)
		return sum_y (vx, 1, 0, m-1, ly, ry);
	int tmx = (tlx + trx) / 2;
	return sum_x (vx*2, tlx, tmx, lx, min(rx,tmx), ly, ry)
		+ sum_x (vx*2+1, tmx+1, trx, max(lx,tmx+1), rx, ly, ry);
}

// EJEMPLO DE UPDATE

// De nuevo, esto es interno, se debería llamar update_x
// Update the Y-axis segment tree for a fixed X-node vx
// vx: fixed X-node index
// lx, rx: X-range covered by vx
// vy: current Y-node index
// ly, ry: Y-range covered by vy
// x, y: coordinates to update
// new_val: new value to set at (x, y)
void update_y (int vx, int lx, int rx, int vy, int ly, int ry, int x, int y, int new_val) {
	if (ly == ry) {
		if (lx == rx)
			t[vx][vy] = new_val;
		else
			t[vx][vy] = t[vx*2][vy] + t[vx*2+1][vy];
	}
	else {
		int my = (ly + ry) / 2;
		if (y <= my)
			update_y (vx, lx, rx, vy*2, ly, my, x, y, new_val);
		else
			update_y (vx, lx, rx, vy*2+1, my+1, ry, x, y, new_val);
		t[vx][vy] = t[vx][vy*2] + t[vx][vy*2+1];
	}
}
 
// O(logn logm + k)
// Update the X-axis segment tree
// vx: current X-node index
// lx, rx: X-range covered by vx
// x, y: coordinates to update
// new_val: new value to set at (x, y)
void update_x (int vx, int lx, int rx, int x, int y, int new_val) {
	if (lx != rx) {
		int mx = (lx + rx) / 2;
		if (x <= mx)
			update_x (vx*2, lx, mx, x, y, new_val);
		else
			update_x (vx*2+1, mx+1, rx, x, y, new_val);
	}
	update_y (vx, lx, rx, 1, 0, m-1, x, y, new_val);
}

\end{minted}

	\needspace{1\baselineskip}
	{
		\phantomsection
		\subsubsection*{Segment Tree}
		\addcontentsline{toc}{subsubsection}{Segment Tree}
	}
	\begin{minted}{cpp}
template <typename T>
struct SegmentTree
{
  // usa, v = 1, tl = 0, tr = n-1 para 0-indexed
  // usa  v = 1, tl = 1, tr = n para 1-indexed pero 
  // asegurate de acceder a nums[tl] correctamente 
  T neuter;
  vector<T> st;
  inline T merge(T a, T b) { return a + b; }
\end{minted}
	\vspace{-12pt}
	\needspace{6\baselineskip}
	\begin{minted}{cpp}
  SegmentTree(vector<T> &nums)
  {
    neuter = 0;
    st.assign(4 * nums.size(), neuter);
    build(1, 0, nums.size() - 1, nums);
  }
\end{minted}
	\vspace{-12pt}
	\needspace{12\baselineskip}
	\begin{minted}{cpp}
  void build(int v, int tl, int tr, vector<T> &nums)
  {
    if (tl == tr)
    {
      st[v] = nums[tl];
      return;
    }
    int tm = tl + (tr - tl) / 2;
    build(v * 2, tl, tm, nums);
    build(v * 2 + 1, tm + 1, tr, nums);
    st[v] = merge(st[2 * v], st[2 * v + 1]);
  }
\end{minted}
	\vspace{-12pt}
	\needspace{12\baselineskip}
	\begin{minted}{cpp}
  void update(int v, int tl, int tr, int pos, T val)
  {
    if (tl == tr)
    {
      st[v] = val;
      return;
    }
    int tm = tl + (tr - tl) / 2;
    if (pos <= tm) update(v * 2, tl, tm, pos, val);
    else update(v * 2 + 1, tm + 1, tr, pos, val);
    st[v] = merge(st[v * 2], st[v * 2 + 1]);
  }
\end{minted}
	\vspace{-12pt}
	\needspace{8\baselineskip}
	\begin{minted}{cpp}
  T query(int v, int tl, int tr, int l, int r)
  {
    if (tl >= l && tr <= r) return st[v];
    if (tl > r || tr < l) return neuter;
    int tm = tl + (tr - tl) / 2;
    return merge(query(v * 2, tl, tm, l, r), 
    query(v * 2 + 1, tm + 1, tr, l, r));
  }
};

\end{minted}

	\needspace{6\baselineskip}
	{
		\phantomsection
		\section*{Graphs}
		\markboth{GRAPHS}{}
		\addcontentsline{toc}{section}{Graphs}
	}
	{
		\phantomsection
		\subsection*{Flows}
		\addcontentsline{toc}{subsection}{Flows}
	}
	{
		\phantomsection
		\subsubsection*{Closure Problem}
		\addcontentsline{toc}{subsubsection}{Closure Problem}
	}
	\begin{minted}{cpp}
/**
 *
 * Problema de cerradura/clausura máxima (o mínima):
 *
 * Una cerradura C de una digráfica D es un conjunto de V(D) tal que no hay
 * flecha que salga de C.
 *
 * El problema de clausura máxima (o análogamente mínima) se trata de encontrar
 * C de peso máximo (o análogamente mínimo), con pesos sobre los vértices.
 *
 * Para pasar de un problema de clausura mínima a clausura máxima, la clausura
 * mínima es equivalente a la clausura máxima de la digráfica conversa (i.e. con
 * las flechas al revés).
 *
 * Para resolver el problema y encontrar el peso de la cerradura máxima se construye
 * la red de flujo con las mismas flechas que D y peso infinito, con los vértices
 * nuevos S y T, que todo vértice v con peso positivo tiene una flecha de S a v
 * con el peso de v, y todo vértice v con peso negativo tiene una flecha de v a
 * T con el valor absoluto del peso de v.
 *
 * El peso de la cerradura máxima es la suma de todos los pesos positivos menos
 * el flujo máximo de la red. Esto corresponde con el corte mínimo necesario
 * del valor máximo, i.e. que es lo menos que podemos "sacrificar" del peso
 * máximo posible.
 *
 * Los vértices en la cerradura máxima son aquellos que están del lado de S en
 * el corte mínimo.
 *
 */
using dig_adj = vector<vector<ll>>;
constexpr ll inf = 1e17;

ll peso_cerradura_max(ll n, dig_adj g, vector<ll> costs)
{
    // Le asignamos a S y T los índices n y n+1
    ll s = n, t = n + 1;
    Dinic flujo(n + 2, s, t);
    // Agrega las aristas en la gráfica
    for (ll v = 0; v < n; v++)
    for (ll u : g[v])
    {
        flujo.add(v, u, inf);
    }
    // Agrega los costos
    for (ll v = 0; v < n; v++)
    {
        ll c = costs[v];
        if (c > 0)
        {
            flujo.add(s, v, c);
        }
        else if (c < 0)
        {
            flujo.add(v, t, -c);
        }
    }

    // Suma de costos
    ll suma = 0;
    for (ll c : costs) if (c > 0) suma += c;

    return suma - flujo.maxFlow();
}

\end{minted}

	\needspace{1\baselineskip}
	{
		\phantomsection
		\subsubsection*{Flow}
		\addcontentsline{toc}{subsubsection}{Flow}
	}
	\begin{minted}{cpp}
#include <bits/stdc++.h>

using namespace std;
typedef long long ll;
typedef vector<int> vi;
typedef pair<int,int> pii;
typedef vector<pii> vii;

#define all(v) v.begin(),v.end()
#define pb(x) push_back(x)

const ll inf = 1e17;

struct flowEdge {
    int u,v;
    ll cap,flow = 0; 
    flowEdge(int u, int v, ll cap) : u(u), v(v), cap(cap) {};
};

struct Dinic{
    vector<flowEdge> edges; 
    vector<vi> adj; 
    int n,s,t; 
    int id = 0;
    vi level, next; 
    
    queue<int> q;
    Dinic(int n, int s, int t) : n(n), s(s), t(t) {
        adj.resize(n);
        level.resize(n);
        next.resize(n);
        fill(all(level),-1); 
        level[s] = 0; 
        q.push(s);
    }

    void add(int u, int v, ll cap){
        edges.emplace_back(u,v,cap);
        edges.emplace_back(v,u,0);
        adj[u].pb(id++); 
        adj[v].pb(id++); 
    }

    bool bfs(){
        while (!q.empty()){
            int curr = q.front();
            q.pop();
            for (auto e : adj[curr]){
                if (edges[e].cap - edges[e].flow  < 1) continue; 
                if (level[edges[e].v] != -1) continue; 
                level[edges[e].v] = level[edges[e].u] + 1; 
                q.push(edges[e].v);
            }
        }
        return level[t] != -1; 
    }

    ll dfs(int u, ll flow){
        if (flow == 0) return 0;
        if (u == t) return flow;
        for (auto& cid = next[u]; cid < adj[u].size(); cid++){ 
            int e = adj[u][cid]; 
            int v = edges[e].v; 

            if (level[edges[e].u] + 1 != level[v] || edges[e].cap - edges[e].flow < 1) continue;
            ll f = dfs(v,min(flow,edges[e].cap - edges[e].flow)); 
            if (f == 0) continue; 
            edges[e].flow += f; 
            edges[e ^ 1].flow -= f;
            return f;
        }
        return 0;
    }

    ll maxFlow(){
        ll flow = 0;
        while (bfs()){
            fill(all(next),0);
            for (ll f = dfs(s,inf); f != 0ll; f = dfs(s,inf)) flow += f;
            fill(all(level),-1); 
            level[s] = 0;
            q.push(s);
        }
        return flow;
    }

};

int main(){
    int n,m; cin >> n >> m;
    Dinic g(n,0,n-1);
    for(int i = 0; i < m ; i++){
        int a,b,c; cin >> a >> b >> c;
        g.add(a, b, c);
    }
    cout << g.maxFlow() << '\n';
}

\end{minted}

	\needspace{1\baselineskip}
	{
		\phantomsection
		\subsubsection*{Min Cost Max Flow (Dinic)}
		\addcontentsline{toc}{subsubsection}{Min Cost Max Flow (Dinic)}
	}
	\begin{minted}{cpp}
// O(Fmlogn) (avg case ~O(sqrt(n)mlogn))
/**
 * Author: Stanford
 * Date: Unknown
 * Source: Stanford Notebook
 * Description: Min-cost max-flow.
 *  If costs can be negative, call setpi before maxflow, but note that negative
 *      cost cycles are not supported.
 *  To obtain the actual flow, look at positive values only.
 * Status: Tested on kattis:mincostmaxflow, stress-tested against another
 *      implementation
 * Time: $O(F E \log(V))$ where F is max flow. $O(VE)$ for setpi.
 */

#define rep(i, a, b) for(int i = a; i < (b); ++i)
#define sz(x) (int)(x).size()

// #include <bits/extc++.h> /// include-line, keep-include

const ll INF = numeric_limits<ll>::max() / 4;

struct MCMF {
	struct edge {
		int from, to, rev;
		ll cap, cost, flow;
	};
	int N;
	vector<vector<edge>> ed;
	vi seen;
	vector<ll> dist, pi;
	vector<edge*> par;

	MCMF(int N) : N(N), ed(N), seen(N), dist(N), pi(N), par(N) {}

	void addEdge(int from, int to, ll cap, ll cost) {
		if (from == to) return;
		ed[from].push_back(edge{ from,to,sz(ed[to]),cap,cost,0 });
		ed[to].push_back(edge{ to,from,sz(ed[from])-1,0,-cost,0 });
	}

	void path(int s) {
		fill(all(seen), 0);
		fill(all(dist), INF);
		dist[s] = 0; ll di;

		__gnu_pbds::priority_queue<pair<ll, int>> q;
		vector<decltype(q)::point_iterator> its(N);
		q.push({ 0, s });

		while (!q.empty()) {
			s = q.top().second; q.pop();
			seen[s] = 1; di = dist[s] + pi[s];
			for (edge& e : ed[s]) if (!seen[e.to]) {
				ll val = di - pi[e.to] + e.cost;
				if (e.cap - e.flow > 0 && val < dist[e.to]) {
					dist[e.to] = val;
					par[e.to] = &e;
					if (its[e.to] == q.end())
						its[e.to] = q.push({ -dist[e.to], e.to });
					else
						q.modify(its[e.to], { -dist[e.to], e.to });
				}
			}
		}
		rep(i,0,N) pi[i] = min(pi[i] + dist[i], INF);
	}

	pair<ll, ll> maxflow(int s, int t) {
		ll totflow = 0, totcost = 0;
		while (path(s), seen[t]) {
			ll fl = INF;
			for (edge* x = par[t]; x; x = par[x->from])
				fl = min(fl, x->cap - x->flow);

			totflow += fl;
			for (edge* x = par[t]; x; x = par[x->from]) {
				x->flow += fl;
				ed[x->to][x->rev].flow -= fl;
			}
		}
		rep(i,0,N) for(edge& e : ed[i]) totcost += e.cost * e.flow;
		return {totflow, totcost/2};
	}

	// If some costs can be negative, call this before maxflow:
	void setpi(int s) { // (otherwise, leave this out)
		fill(all(pi), INF); pi[s] = 0;
		int it = N, ch = 1; ll v;
		while (ch-- && it--)
			rep(i,0,N) if (pi[i] != INF)
			  for (edge& e : ed[i]) if (e.cap)
				  if ((v = pi[i] + e.cost) < pi[e.to])
					  pi[e.to] = v, ch = 1;
		assert(it >= 0); // negative cost cycle
	}
};

\end{minted}

	\needspace{1\baselineskip}
	{
		\phantomsection
		\subsubsection*{Min Cost Max Flow (Edmond Karps)}
		\addcontentsline{toc}{subsubsection}{Min Cost Max Flow (Edmond Karps)}
	}
	\begin{minted}{cpp}
// O(Fnm)
struct Edge
{
    int from, to, capacity, cost;
};

vector<vector<int>> adj, cost, capacity;

const int INF = 1e9;

void shortest_paths(int n, int v0, vector<int>& d, vector<int>& p) {
    d.assign(n, INF);
    d[v0] = 0;
    vector<bool> inq(n, false);
    queue<int> q;
    q.push(v0);
    p.assign(n, -1);

    while (!q.empty()) {
        int u = q.front();
        q.pop();
        inq[u] = false;
        for (int v : adj[u]) {
            if (capacity[u][v] > 0 && d[v] > d[u] + cost[u][v]) {
                d[v] = d[u] + cost[u][v];
                p[v] = u;
                if (!inq[v]) {
                    inq[v] = true;
                    q.push(v);
                }
            }
        }
    }
}

int min_cost_flow(int N, vector<Edge> edges, int K, int s, int t) {
    adj.assign(N, vector<int>());
    cost.assign(N, vector<int>(N, 0));
    capacity.assign(N, vector<int>(N, 0));
    for (Edge e : edges) {
        adj[e.from].push_back(e.to);
        adj[e.to].push_back(e.from);
        cost[e.from][e.to] = e.cost;
        cost[e.to][e.from] = -e.cost;
        capacity[e.from][e.to] = e.capacity;
    }

    int flow = 0;
    int cost = 0;
    vector<int> d, p;
    while (flow < K) {
        shortest_paths(N, s, d, p);
        if (d[t] == INF)
            break;

        // find max flow on that path
        int f = K - flow;
        int cur = t;
        while (cur != s) {
            f = min(f, capacity[p[cur]][cur]);
            cur = p[cur];
        }

        // apply flow
        flow += f;
        cost += f * d[t];
        cur = t;
        while (cur != s) {
            capacity[p[cur]][cur] -= f;
            capacity[cur][p[cur]] += f;
            cur = p[cur];
        }
    }

    if (flow < K)
        return -1;
    else
        return cost;
}

\end{minted}

	\needspace{1\baselineskip}
	{
		\phantomsection
		\subsubsection*{Min Cut}
		\addcontentsline{toc}{subsubsection}{Min Cut}
	}
	\begin{minted}{cpp}
/**
 * Author: Simon Lindholm
 * Date: 2021-01-09
 * License: CC0
 * Source: https://en.wikipedia.org/wiki/Stoer%E2%80%93Wagner_algorithm
 * Description: Find a global minimum cut in an undirected graph, as
 *              represented by an adjacency matrix.
 * Time: O(V^3)
 * Status: Stress-tested together with GomoryHu
 */
#pragma once

pair<int, vi> globalMinCut(vector<vi> mat) {
	pair<int, vi> best = {INT_MAX, {}};
	int n = sz(mat);
	vector<vi> co(n);
	rep(i,0,n) co[i] = {i};
	rep(ph,1,n) {
		vi w = mat[0];
		size_t s = 0, t = 0;
		rep(it,0,n-ph) { // O(V^2) -> O(E log V) with prio. queue
			w[t] = INT_MIN;
			s = t, t = max_element(all(w)) - w.begin();
			rep(i,0,n) w[i] += mat[t][i];
		}
		best = min(best, {w[t] - mat[t][t], co[t]});
		co[s].insert(co[s].end(), all(co[t]));
		rep(i,0,n) mat[s][i] += mat[t][i];
		rep(i,0,n) mat[i][s] = mat[s][i];
		mat[0][t] = INT_MIN;
	}
	return best;
}

\end{minted}

	\needspace{1\baselineskip}
	{
		\phantomsection
		\subsubsection*{Push Relabel Flow}
		\addcontentsline{toc}{subsubsection}{Push Relabel Flow}
	}
	\begin{minted}{cpp}
/**
 * Author: Simon Lindholm
 * Date: 2015-02-24
 * License: CC0
 * Source: Wikipedia, tinyKACTL
 * Description: Push-relabel using the highest label selection rule and the gap
 *     heuristic. Quite fast in practice. To obtain the actual flow, look at
 *     positive values only.
 * Time: $O(V^2\sqrt E)$
 * Status: Tested on Kattis and SPOJ, and stress-tested
 */

#define rep(i, a, b) for(int i = a; i < (b); ++i)
#define sz(x) (int)(x).size()

struct PushRelabel {
	struct Edge {
		int dest, back;
		ll f, c;
	};
	vector<vector<Edge>> g;
	vector<ll> ec;
	vector<Edge*> cur;
	vector<vi> hs; vi H;
	PushRelabel(int n) : g(n), ec(n), cur(n), hs(2*n), H(n) {}

	void addEdge(int s, int t, ll cap, ll rcap=0) {
		if (s == t) return;
		g[s].push_back({t, sz(g[t]), 0, cap});
		g[t].push_back({s, sz(g[s])-1, 0, rcap});
	}

	void addFlow(Edge& e, ll f) {
		Edge &back = g[e.dest][e.back];
		if (!ec[e.dest] && f) hs[H[e.dest]].push_back(e.dest);
		e.f += f; e.c -= f; ec[e.dest] += f;
		back.f -= f; back.c += f; ec[back.dest] -= f;
	}
	ll calc(int s, int t) {
		int v = sz(g); H[s] = v; ec[t] = 1;
		vi co(2*v); co[0] = v-1;
		rep(i,0,v) cur[i] = g[i].data();
		for (Edge& e : g[s]) addFlow(e, e.c);

		for (int hi = 0;;) {
			while (hs[hi].empty()) if (!hi--) return -ec[s];
			int u = hs[hi].back(); hs[hi].pop_back();
			while (ec[u] > 0)  // discharge u
				if (cur[u] == g[u].data() + sz(g[u])) {
					H[u] = 1e9;
					for (Edge& e : g[u]) if (e.c && H[u] > H[e.dest]+1)
						H[u] = H[e.dest]+1, cur[u] = &e;
					if (++co[H[u]], !--co[hi] && hi < v)
						rep(i,0,v) if (hi < H[i] && H[i] < v)
							--co[H[i]], H[i] = v + 1;
					hi = H[u];
				} else if (cur[u]->c && H[u] == H[cur[u]->dest]+1)
					addFlow(*cur[u], min(ec[u], cur[u]->c));
				else ++cur[u];
		}
	}
	bool leftOfMinCut(int a) { return H[a] >= sz(g); }
};

\end{minted}

	\needspace{3\baselineskip}
	{
		\phantomsection
		\subsection*{Graphs algorthms}
		\addcontentsline{toc}{subsection}{Graphs algorthms}
	}
	{
		\phantomsection
		\subsubsection*{Acyclic Ordering}
		\addcontentsline{toc}{subsubsection}{Acyclic Ordering}
	}
	\begin{minted}{cpp}
using ll = long long;
using dig_adj = vector<vector<ll>>;
using vll = vector<ll>;

void dfsa_aux(dig_adj &d, vll &pred, vll &tvisit, vll &texpl, vll &verts, ll &time, ll &i, ll v);
/**
 * Algoritmo de Knuth para ordenamiento acíclico. O(n+m). Variación de DFS.
 *
 * Sea D una digráfica y sea x_1, x_2, ..., x_n un ordenamiento de sus vértices.
 * Decimos que el ordenamiento es acíclico si para cada flecha x_ix_j en D, i<j.
 * Existe si y sólo si la digráfica no tiene un ciclo dirigido.
 *
 * Como toda digráfica acíclica tiene fuente y sumidero, la idea del algoritmo
 * es escoger fuentes, "quitarlas" y repetir.
 *
 * Como este algoritmo es correcto si la digráfica es acíclica, se puede
 * verificar si la digráfica es acíclica al ejecutarlo y revisar si el
 * ordenamiento es acíclico.
 *
 * Este algoritmo es auxiliar al algoritmo de componentes fuertemente conexas.
 */
vll dfsa(ll n, dig_adj &d)
{
    // pred - predecesor || tvisit - tiempo de visita || texpl - tiempo de exploración
    vll pred(n,-1), tvisit(n,0), texpl(n,0), verts(n,-1);
    ll time = 0, i = n;

    for (ll v = 0; v < n; v++)
    {
        if (tvisit[v]) continue;
        dfsa_aux(d, pred, tvisit, texpl, verts, time, i, v);
    }

    /**
     * Hasta aquí tenemos:
     *   lista de predecesores - pred
     *   tiempos de visita de cada vértice - tvisit
     *   tiempo en el que se explora (tras sus hijos) el vértice - texpl
     *   el ordenamiento acíclico - verts
     */
    return verts;
}
// Auxiliar que realiza la búsqueda por profundidad recursivamente
void dfsa_aux(dig_adj &d, vll &pred, vll &tvisit, vll &texpl, vll &verts, ll &time, ll &i, ll v)
{
    tvisit[v] = ++time;
    for (ll u : d[v])
    {
        if (tvisit[u]) continue;
        pred[u] = v;
        dfsa_aux(d, pred, tvisit, texpl, verts, time, i, u);
    }
    texpl[v] = ++time;
    verts[--i] = v;       
}

\end{minted}

	\needspace{1\baselineskip}
	{
		\phantomsection
		\subsubsection*{Bellman-Ford}
		\addcontentsline{toc}{subsubsection}{Bellman-Ford}
	}
	\begin{minted}{cpp}
struct Edge {
    int a, b, cost;
};

int n, m, v; // #verts, #edges, vert ini
vector<Edge> edges;
const int INF = 1000000000;

// Bellman Ford que sólo nos dice el costo desde v a los demás
void solve()
{
    vector<int> d(n, INF);
    d[v] = 0;
    for (;;) {
        bool any = false;

        for (Edge e : edges)
            if (d[e.a] < INF)
                if (d[e.b] > d[e.a] + e.cost) {
                    d[e.b] = d[e.a] + e.cost;
                    any = true;
                }

        if (!any)
            break;
    }
    // display d, for example, on the screen
}

// Bellman Ford con la ruta desde v hacia t
void solve()
{
    vector<int> d(n, INF);
    d[v] = 0;
    vector<int> p(n, -1);

    for (;;) {
        bool any = false;
        for (Edge e : edges)
            if (d[e.a] < INF)
                if (d[e.b] > d[e.a] + e.cost) {
                    d[e.b] = d[e.a] + e.cost;
                    p[e.b] = e.a;
                    any = true;
                }
        if (!any)
            break;
    }

    if (d[t] == INF)
        cout << "No path from " << v << " to " << t << ".";
    else {
        vector<int> path;
        for (int cur = t; cur != -1; cur = p[cur])
            path.push_back(cur);
        reverse(path.begin(), path.end());

        cout << "Path from " << v << " to " << t << ": ";
        for (int u : path)
            cout << u << ' ';
    }
}

// Bellman Ford con detección de ciclos negativos
void solve()
{
    vector<int> d(n, INF);
    d[v] = 0;
    vector<int> p(n, -1);
    int x;
    for (int i = 0; i < n; ++i) {
        x = -1;
        for (Edge e : edges)
            if (d[e.a] < INF)
                if (d[e.b] > d[e.a] + e.cost) {
                    d[e.b] = max(-INF, d[e.a] + e.cost);
                    p[e.b] = e.a;
                    x = e.b;
                }
    }

    if (x == -1)
        cout << "No negative cycle from " << v;
    else {
        int y = x;
        for (int i = 0; i < n; ++i)
            y = p[y];

        vector<int> path;
        for (int cur = y;; cur = p[cur]) {
            path.push_back(cur);
            if (cur == y && path.size() > 1)
                break;
        }
        reverse(path.begin(), path.end());

        cout << "Negative cycle: ";
        for (int u : path)
            cout << u << ' ';
    }
}

/////////////////////////////

// Shortest path faster algorithm (promedio O(m), peor caso O(nm))
const int INF = 1000000000;
vector<vector<pair<int, int>>> adj;
bool spfa(int s, vector<int>& d) {
    int n = adj.size();
    d.assign(n, INF);
    vector<int> cnt(n, 0);
    vector<bool> inqueue(n, false);
    queue<int> q;

    d[s] = 0;
    q.push(s);
    inqueue[s] = true;
    while (!q.empty()) {
        int v = q.front();
        q.pop();
        inqueue[v] = false;

        for (auto edge : adj[v]) {
            int to = edge.first;
            int len = edge.second;

            if (d[v] + len < d[to]) {
                d[to] = d[v] + len;
                if (!inqueue[to]) {
                    q.push(to);
                    inqueue[to] = true;
                    cnt[to]++;
                    if (cnt[to] > n)
                        return false;  // negative cycle
                }
            }
        }
    }
    return true;
}

\end{minted}

	\needspace{1\baselineskip}
	{
		\phantomsection
		\subsubsection*{Cycle Detection Directed}
		\addcontentsline{toc}{subsubsection}{Cycle Detection Directed}
	}
	\begin{minted}{cpp}
int n;
vector<vector<int>> adj;
vector<char> color;
vector<int> parent;
int cycle_start, cycle_end;

bool dfs(int v) {
    color[v] = 1;
    for (int u : adj[v]) {
        if (color[u] == 0) {
            parent[u] = v;
            if (dfs(u))
                return true;
        } else if (color[u] == 1) {
            cycle_end = v;
            cycle_start = u;
            return true;
        }
    }
    color[v] = 2;
    return false;
}

void find_cycle() {
    color.assign(n, 0);
    parent.assign(n, -1);
    cycle_start = -1;

    for (int v = 0; v < n; v++) {
        if (color[v] == 0 && dfs(v))
            break;
    }

    if (cycle_start == -1) {
        cout << "Acyclic" << endl;
    } else {
        vector<int> cycle;
        cycle.push_back(cycle_start);
        for (int v = cycle_end; v != cycle_start; v = parent[v])
            cycle.push_back(v);
        cycle.push_back(cycle_start);
        reverse(cycle.begin(), cycle.end());

        cout << "Cycle found: ";
        for (int v : cycle)
            cout << v << " ";
        cout << endl;
    }
}
\end{minted}

	\needspace{1\baselineskip}
	{
		\phantomsection
		\subsubsection*{Cycle Detection Undirected}
		\addcontentsline{toc}{subsubsection}{Cycle Detection Undirected}
	}
	\begin{minted}{cpp}
int n;
vector<vector<int>> adj;
vector<bool> visited;
vector<int> parent;
int cycle_start, cycle_end;

bool dfs(int v, int par) { // passing vertex and its parent vertex
    visited[v] = true;
    for (int u : adj[v]) {
        if(u == par) continue; // skipping edge to parent vertex
        if (visited[u]) {
            cycle_end = v;
            cycle_start = u;
            return true;
        }
        parent[u] = v;
        if (dfs(u, parent[u]))
            return true;
    }
    return false;
}

void find_cycle() {
    visited.assign(n, false);
    parent.assign(n, -1);
    cycle_start = -1;

    for (int v = 0; v < n; v++) {
        if (!visited[v] && dfs(v, parent[v]))
            break;
    }

    if (cycle_start == -1) {
        cout << "Acyclic" << endl;
    } else {
        vector<int> cycle;
        cycle.push_back(cycle_start);
        for (int v = cycle_end; v != cycle_start; v = parent[v])
            cycle.push_back(v);
        cycle.push_back(cycle_start);

        cout << "Cycle found: ";
        for (int v : cycle)
            cout << v << " ";
        cout << endl;
    }
}
\end{minted}

	\needspace{1\baselineskip}
	{
		\phantomsection
		\subsubsection*{Dijkstra}
		\addcontentsline{toc}{subsubsection}{Dijkstra}
	}
	\begin{minted}{cpp}
const int INF = 1000000000;
vector<vector<pair<int, int>>> adj;

void dijkstra(int s, vector<int> & d, vector<int> & p) {
    int n = adj.size();
    d.assign(n, INF);
    p.assign(n, -1);

    d[s] = 0;
    using pii = pair<int, int>;
    priority_queue<pii, vector<pii>, greater<pii>> q;
    q.push({0, s});
    while (!q.empty()) {
        int v = q.top().second;
        int d_v = q.top().first;
        q.pop();
        if (d_v != d[v])
            continue;

        for (auto edge : adj[v]) {
            int to = edge.first;
            int len = edge.second;

            if (d[v] + len < d[to]) {
                d[to] = d[v] + len;
                p[to] = v;
                q.push({d[to], to});
            }
        }
    }
}
\end{minted}

	\needspace{1\baselineskip}
	{
		\phantomsection
		\subsubsection*{Euler tour}
		\addcontentsline{toc}{subsubsection}{Euler tour}
	}
	\begin{minted}{cpp}
const int N = 2e5 + 1;
int timer = 0;
int start[N];
int endd[N];
ll arreglo_segtree[N];
ll values[N];
vector<int> graph[N];
void euler_tour(int u, int p)
{
  start[u] = timer++;
  for (int v : graph[u])
    if (v != p)
      euler_tour(v, u);
  endd[u] = timer;
}

int main()
{
  int n, q; cin >> n >> q;
  for (int i = 0; i < n; i++)
    cin >> values[i];
  
  // Gŕafica con nodos 0-indexados
  for (int i = 0; i < n - 1; i++)
  {
    int u, v; cin >> u >> v;
    u--;
    v--;
    graph[u].push_back(v);
    graph[v].push_back(u);
  }

  // Empezamos el euler tour con la raiz
  int root = 0;
  euler_tour(root, -1);
  // Para este punto start[i] y end[i] tienen el tiempo en el que se visito y se terminó de visitar el vertice i en el árbol 
  // end[i] - start[i] = Tamaaño del subarbol de i (incluyendolo)

  // Construimos un arreglo con los valores de los nodos, con el orden dado por start, que hará que los nodos estén contiguos y podamos hacer uso de un SegmentTree
  for (int i = 0; i < n; i++)
    arreglo_segtree[start[i]] = values[i];

  SegmentTree<ll> st(arreglo_segtree, n);
  // start[i] = indice desde donde comienza el subarbol de i (incluyendolo) en arreglo_segtree
  // endd[i] - 1 = incide donde termina el subarbol de i en arreglo_segtree (es importante restarle uno)

  while (q--)
  {
    int op; cin >> op;
    int node; cin >> node;
    node--;

    int l = start[node]; 
    int r = endd[node] - 1;

    if (op == 1)
    {
      ll val;
      cin >> val;
      // l es la posición donde está el nodo en arreglo_segtree
      st.update(1, 0, n - 1, l, val);
    }
    else
    {
      // En caso de que no se incluya a nodo en el subarbol
      // puede pasar que el subarbol sea vacío
      if (r < l)
      {
        cout << 0 << endl;
        continue;
      }
      cout << st.query(1, 0, n - 1, l, r) << endl;
    }
  }
  return 0;
}
\end{minted}

	\needspace{1\baselineskip}
	{
		\phantomsection
		\subsubsection*{Floyd Warshall}
		\addcontentsline{toc}{subsubsection}{Floyd Warshall}
	}
	\begin{minted}{cpp}
const int INF = 1e9;
int n; 

vector<vector<int>> d(n, vector<int>(n, INF));
// Inicializamos la matriz de distancias:
// - d[i][j] = INF si no hay un camino directo conocido de i a j
// - d[i][i] = 0 para todo i (la distancia de un nodo a sí mismo es 0)
// - d[i][j] = peso de la arista de i a j, si existe una arista directa

// Floyd-Warshall: encuentra distancias mínimas entre todos los pares (i, j)
for (int k = 0; k < n; ++k) {
    // Usamos el nodo k como intermediario
    for (int i = 0; i < n; ++i) {
        for (int j = 0; j < n; ++j) {
            // Solo procesamos si los caminos i -> k y k -> j existen
            if (d[i][k] < INF && d[k][j] < INF)
                // Verificamos si ir de i a j pasando por k es más barato
                d[i][j] = min(d[i][j], d[i][k] + d[k][j]);
        }
    }
}

\end{minted}

	\needspace{1\baselineskip}
	{
		\phantomsection
		\subsubsection*{Fuertemente Conexa Kosaraju-Sharir}
		\addcontentsline{toc}{subsubsection}{Fuertemente Conexa Kosaraju-Sharir}
	}
	\begin{minted}{cpp}
// NOTE: Hace uso de Acyclic Ordering (dfsa)

using ll = long long;
using dig_adj = vector<vector<ll>>;
using vll = vector<ll>;

void sca_aux(dig_adj &conv, vll &pred, vll &tvisit, vll &texpl, ll &time, vll &componente, ll comp_idx, ll v);
/**
 * Componentes fuertemente conexas (Kosaraju-Sharir). O(n+m). Variación de DFS.
 *
 * Además, construimos la gráfica de bloques fuertes pero se indica en un
 * comentario donde podemos deternernos si sólo queremos las componentes.
 * A veces a esto se le llama la condensación de la digráfica.
 *
 * Aquí estoy realizando una pequeña modificación para el problema: construyo
 * la digráfica de bloques fuertes con el costo acumulado de todos los vértices
 * en cada bloque (pues todo bloque únicamente puede ser elegido todo junto).
 *
 * Recuerde que la digráfica conversa es la que tiene las flechas invertidas
 */
dig_adj sca(ll n, dig_adj &d, dig_adj &d_converse)
{
    vll ord_aciclico = dfsa(n, d);

    // pred - predecesor || tvisit - tiempo de visita || texpl - tiempo de exploración
    vll pred(n,-1), tvisit(n,0), texpl(n,0);
    vll componente(n,-1); // a cuál componente pertenece cada vértice
    ll time = 0, componente_idx = 0;

    // La consulta de los vértices debe ser sobre el orden acíclico
    for (ll v : ord_aciclico)
    {
        if (tvisit[v]) continue;
        sca_aux(d_converse, pred, tvisit, texpl, time, componente, componente_idx++, v);
    }

    /**
     * Hasta aquí tenemos:
     *   lista de predecesores - pred
     *   tiempos de visita de cada vértice - tvisit
     *   tiempo en el que se explora (tras sus hijos) el vértice - texpl
     *   a qué componente pertenece cada vértice - componente
     *   cuántas componentes hay - componente_idx
     */

    dig_adj bloques(componente_idx);
    // Revisamos todas las aristas y agregamos las que unen componentes
    for (ll u = 0; u < n; u++)
    for (ll v : d[u])
    {
        ll cu = componente[u];
        ll cv = componente[v];
        if (cu == cv) continue;
        bloques.m++;
        bloques[cu].push_back(cv);
    }

    return bloques;
}
// Auxiliar que realiza la búsqueda por profundidad recursivamente
void sca_aux(dig_adj &conv, vll &pred, vll &tvisit, vll &texpl, ll &time, vll &componente, ll comp_idx, ll v)
{
    tvisit[v] = ++time;
    componente[v] = comp_idx;
    for (ll u : conv[v])
    {
        if (tvisit[u]) continue;
        pred[u] = v;
        sca_aux(conv, pred, tvisit, texpl, time, componente, comp_idx, u);
    }
    texpl[v] = ++time;
}

\end{minted}

	\needspace{1\baselineskip}
	{
		\phantomsection
		\subsubsection*{Havel-hakin}
		\addcontentsline{toc}{subsubsection}{Havel-hakin}
	}
	\begin{minted}{py}
import sys

def todos_son_ceros(t):
    for d in t:
        if d != 0:
            return False
    return True

def esGrafica(grados):
    if sum(grados) % 2 != 0:
        return False
    
    grados = sorted(grados, reverse=True)    
    
    if todos_son_ceros(grados):
        return True

    d = grados[0]
    restantes = list(grados[1:])

    if d > len(restantes):
        return False
    
    for i in range(d):
        restantes[i] -= 1
        if restantes[i] < 0:
            return False
        
    return esGrafica(restantes)

def main():
    data = sys.stdin.read().strip().split()
    grad = list(map(int, data))
    print(esGrafica(grad))

if __name__ == "__main__":
    main()

\end{minted}

	\needspace{1\baselineskip}
	{
		\phantomsection
		\subsubsection*{Is Bipartite}
		\addcontentsline{toc}{subsubsection}{Is Bipartite}
	}
	\begin{minted}{cpp}
int n; 
vector<vector<int>> adj; 

vector<int> side(n, -1); // side[i] = 0 o 1 si ya fue asignado, -1 si no ha sido visitado
bool is_bipartite = true; 

queue<int> q; 
// Recorremos todos los componentes del grafo
for (int st = 0; st < n; ++st) {
    if (side[st] == -1) {
        // Si el nodo aún no ha sido visitado, iniciamos BFS desde él
        q.push(st);
        side[st] = 0; // Le asignamos el color 0

        while (!q.empty()) {
            int v = q.front(); q.pop();
            for (int u : adj[v]) {
                if (side[u] == -1) {
                    // Si el vecino no ha sido visitado, lo coloreamos con el color opuesto
                    side[u] = side[v] ^ 1; // Alternamos entre 0 y 1
                    q.push(u);
                } else {
                    // Si el vecino ya está coloreado y tiene el mismo color, no es bipartito
                    if (side[u] == side[v])
                        is_bipartite = false;
                }
            }
        }
    }
}
cout << (is_bipartite ? "YES" : "NO") << endl; 

\end{minted}

	\needspace{1\baselineskip}
	{
		\phantomsection
		\subsubsection*{Tarjan}
		\addcontentsline{toc}{subsubsection}{Tarjan}
	}
	\begin{minted}{cpp}
#include <iostream>
#include <vector>
#include <set>
using namespace std;

const int MAXN = 100005;  // ajusta según el tamaño de tu grafo
vector<int> grafo[MAXN];
bool visitado[MAXN];
int tin[MAXN], low[MAXN];
int timer = 0;
set<int> puntos_articulacion;  // para evitar repetidos

void dfs(int u, int padre = -1) {
    visitado[u] = true;
    tin[u] = low[u] = timer++;
    int hijos = 0;

    for (int v : grafo[u]) {
        if (v == padre) continue;

        if (visitado[v]) {
            // v ya fue visitado: actualiza low
            low[u] = min(low[u], tin[v]);
        } else {
            dfs(v, u);
            low[u] = min(low[u], low[v]);
            if (low[v] >= tin[u] && padre != -1) {
                puntos_articulacion.insert(u);
            }
            ++hijos;
        }
    }

    // Caso especial: si u es la raíz y tiene más de un hijo
    if (padre == -1 && hijos > 1) {
        puntos_articulacion.insert(u);
    }
}

int main() {
    int n, m;
    cin >> n >> m; // número de nodos y aristas
    for (int i = 0; i < m; ++i) {
        int u, v;
        cin >> u >> v;
        grafo[u].push_back(v);
        grafo[v].push_back(u);  
    }

    for (int i = 0; i < n; ++i) {
        if (!visitado[i]) dfs(i);
    }

    cout << "Puntos de articulación:\n";
    for (int p : puntos_articulacion) {
        cout << p << "\n";
    }

    return 0;
}

\end{minted}

	\needspace{1\baselineskip}
	{
		\phantomsection
		\subsubsection*{Tarjan puentes}
		\addcontentsline{toc}{subsubsection}{Tarjan puentes}
	}
	\begin{minted}{cpp}
#include <iostream>
#include <vector>
using namespace std;

const int MAXN = 100005;
vector<int> grafo[MAXN];
bool visitado[MAXN];
int tin[MAXN], low[MAXN];
int timer = 0;
vector<pair<int, int>> puentes;

void dfs(int u, int padre = -1) {
    visitado[u] = true;
    tin[u] = low[u] = timer++;

    for (int v : grafo[u]) {
        if (v == padre) continue;

        if (visitado[v]) {
            // back edge
            low[u] = min(low[u], tin[v]);
        } else {
            dfs(v, u);
            low[u] = min(low[u], low[v]);

            if (low[v] > tin[u]) {
                // u-v es un puente
                puentes.push_back({u, v});
            }
        }
    }
}

int main() {
    int n, m;
    cin >> n >> m; // número de nodos y aristas

    for (int i = 0; i < m; ++i) {
        int u, v;
        cin >> u >> v;
        grafo[u].push_back(v);
        grafo[v].push_back(u); // no dirigido
    }

    for (int i = 0; i < n; ++i) {
        if (!visitado[i]) dfs(i);
    }

    cout << "Puentes encontrados:\n";
    for (auto [u, v] : puentes) {
        cout << u << " - " << v << "\n";
    }

    return 0;
}

\end{minted}

	\needspace{1\baselineskip}
	{
		\phantomsection
		\subsubsection*{Topological Sort}
		\addcontentsline{toc}{subsubsection}{Topological Sort}
	}
	\begin{minted}{cpp}
int n; // number of vertices
vector<vector<int>> adj; // adjacency list of graph
vector<bool> visited;
vector<int> ans;

void dfs(int v) {
    visited[v] = true;
    for (int u : adj[v]) {
        if (!visited[u]) {
            dfs(u);
        }
    }
    ans.push_back(v);
}

// Se asume que la gráfica es acíclica 
void topological_sort() {
    visited.assign(n, false);
    ans.clear();
    for (int i = 0; i < n; ++i) {
        if (!visited[i]) {
            dfs(i);
        }
    }
    reverse(ans.begin(), ans.end());
}
\end{minted}

	\needspace{1\baselineskip}
	{
		\phantomsection
		\subsubsection*{Weighted Bipartite Matching (Hungarian)}
		\addcontentsline{toc}{subsubsection}{Weighted Bipartite Matching (Hungarian)}
	}
	\begin{minted}{cpp}
/**
 * Author: Benjamin Qi, chilli
 * Date: 2020-04-04
 * License: CC0
 * Source: https://github.com/bqi343/USACO
 * Description: Given a weighted bipartite graph, matches every node on
 * the left with a node on the right such that no
 * nodes are in two matchings and the sum of the edge weights is minimal. Takes
 * cost[N][M], where cost[i][j] = cost for L[i] to be matched with R[j] and
 * returns (min cost, match), where L[i] is matched with
 * R[match[i]]. Negate costs for max cost. Requires $N \le M$.
 * Time: O(N^2M)
 * Status: Tested on kattis:cordonbleu, stress-tested
 */
#define rep(i, a, b) for(int i = a; i < (b); ++i)
#define sz(x) (int)(x).size()

pair<int, vi> hungarian(const vector<vi> &a) {
	if (a.empty()) return {0, {}};
	int n = sz(a) + 1, m = sz(a[0]) + 1;
	vi u(n), v(m), p(m), ans(n - 1);
	rep(i,1,n) {
		p[0] = i;
		int j0 = 0; // add "dummy" worker 0
		vi dist(m, INT_MAX), pre(m, -1);
		vector<bool> done(m + 1);
		do { // dijkstra
			done[j0] = true;
			int i0 = p[j0], j1, delta = INT_MAX;
			rep(j,1,m) if (!done[j]) {
				auto cur = a[i0 - 1][j - 1] - u[i0] - v[j];
				if (cur < dist[j]) dist[j] = cur, pre[j] = j0;
				if (dist[j] < delta) delta = dist[j], j1 = j;
			}
			rep(j,0,m) {
				if (done[j]) u[p[j]] += delta, v[j] -= delta;
				else dist[j] -= delta;
			}
			j0 = j1;
		} while (p[j0]);
		while (j0) { // update alternating path
			int j1 = pre[j0];
			p[j0] = p[j1], j0 = j1;
		}
	}
	rep(j,1,m) if (p[j]) ans[p[j] - 1] = j - 1;
	return {-v[0], ans}; // min cost
}

\end{minted}

	\needspace{3\baselineskip}
	{
		\phantomsection
		\subsection*{Trees}
		\addcontentsline{toc}{subsection}{Trees}
	}
	{
		\phantomsection
		\subsubsection*{Centroid Decomposition}
		\addcontentsline{toc}{subsubsection}{Centroid Decomposition}
	}
	\begin{minted}{cpp}
// A partir de un árbol como gráfica, construye en O(nlogn) un árbol enraizado
// balanceado (de altura O(logn)) dividiéndolo recursivamente en centroides

using graf_t = vector<vector<int>>;

struct ctree_n
{
    int id, parent;
    vector<int> children;
    ctree_n(int id_, int parent_ = -1) : id(id_), parent(parent_), children({}) {}
};

struct ctree
{
    vector<int> size; // size of node[i] subtree
    vector<ctree_n> node;
    int root;

    ctree(const graf_t& graf)
        : size(graf.size(),0)
    {
        int n = graf.size();
        node.reserve(n);
        vector<bool> visited(n,false);
        for (int i = 0; i < n; i++) node.emplace_back(i);
        decompose(0, -1, visited, graf);
    }

private:
    void decompose(int u, int parent, vector<bool>& visited, const graf_t& graf)
    {
        vector<bool> viscp1(visited), viscp2(visited);
        int total = subtree_size(u, -1, viscp1, graf);
        int centroid = find_centroid(u, -1, total, viscp2, graf);
        visited[centroid] = true;
        node[centroid].parent = parent;
        if (parent == -1) node[centroid].parent = root = centroid;
        else node[parent].children.push_back(centroid);
        for (int v : graf[centroid])
            if (!visited[v]) decompose(v, centroid, visited, graf);
    }

    int subtree_size(int u, int p, vector<bool>& visited, const graf_t& graf)
    {
        size[u] = 1;
        for (int v : graf[u])
            if (v != p && !visited[v])
                size[u] += subtree_size(v, u, visited, graf);
        return size[u];
    }

    int find_centroid(int u, int p, int total, vector<bool>& visited, const graf_t& graf)
    {
        for (int v : graf[u])
            if (v != p and !visited[v] and size[v] > total/2)
                return find_centroid(v, u, total, visited, graf);
        return u;
    }
};

\end{minted}

	\needspace{1\baselineskip}
	{
		\phantomsection
		\subsubsection*{Farthest node for each node}
		\addcontentsline{toc}{subsubsection}{Farthest node for each node}
	}
	\begin{minted}{cpp}
// dist[0][i] = distancia del nodo 'a' hacia todos los nodos i
// dist[1][i] = distancia del nodo 'b'' hacai todos los nodos i
int dist[2][N];
vector<int> graph[N];

// f_n = nodo más lejano
int dfs(int u, int p, int d, int i)
{
  dist[i][u] = d;
  int f_n = -1;
  for (int v : graph[u])
  {
    if (v != p)
    {
      int x = dfs(v, u, d + 1, i);
      if (f_n == -1 || dist[i][x] > dist[i][f_n]) f_n = x;
    }
  }
  return f_n == -1 ? u : f_n;
}

void calculate_diameter()
{
  int n; cin >> n;
  for (int i = 0; i < n - 1; i++)
  {
    int a, b;
    cin >> a >> b;
    graph[a-1].push_back(b-1);
    graph[b-1].push_back(a-1);
  }
  // Encontramos el nodo 'a' haciendo un dfs desde un nodo arbitrario
  int a = dfs(0, 0, 0, 0);

  // Encontramos el nodo 'b' haciendo un dfs desde 'a', y su vez calculamos la distancia de 'a' hacia todos los demás nodos
  int b = dfs(a, a, 0, 0);

  // Finalmente calculamos la distncia del nodo 'b' hacia todos los demas
  dfs(b, b, 0, 1);

  // dist(a, b) = Longitud del diametro del árbol
  // max(dist[0][u], dist[1][u]) = distancia de u hacia su nodo más lejano
}

\end{minted}

	\needspace{1\baselineskip}
	{
		\phantomsection
		\subsubsection*{Heavy light decomposition}
		\addcontentsline{toc}{subsubsection}{Heavy light decomposition}
	}
	\begin{minted}{cpp}
// En un árbol enraizado, marcamos a una única arista hacia un hijo como pesada
// si cuenta con la mayor cantidad de vértices en su subárbol. Esto nos forma
// trayectorias pesadas, sobre las cuales podemos agregar segment-trees de
// prefijos para hacer updates/queries de trayectorias, pues se puede mostrar
// que cualquier trayectoria cuenta con a lo más O(logn) de trayectorias
// pesadas. Las aristas no-pesadas les llamamos ligeras.
// La descomposición toma O(n).

// En varios casos se puede usar un segment tree sobre el euler tour en lugar
// de hacer esto.

// parent[v] es el padre de v, depth[v] es la profundidad en v, heavy[v] indica
// al hijo con el cual <v,heavy[v]> es arista pesada, head[v] es donde termina
// el heavy path de v, pos[v] es la posición en el segment-tree de dicho heavy
// path completo. En decompose se llena head y pos.
vector<int> parent, depth, heavy, head, pos;
int cur_pos;

// esto se llama desde init, no directamente
int dfs(int v, vector<vector<int>> const& adj) {
    int size = 1;
    int max_c_size = 0;
    for (int c : adj[v]) {
        if (c != parent[v]) {
            parent[c] = v, depth[c] = depth[v] + 1;
            int c_size = dfs(c, adj);
            size += c_size;
            if (c_size > max_c_size)
                max_c_size = c_size, heavy[v] = c;
        }
    }
    return size;
}

// esto se llama desde init, no directamente
void decompose(int v, int h, vector<vector<int>> const& adj) {
    head[v] = h, pos[v] = cur_pos++;
    if (heavy[v] != -1)
        decompose(heavy[v], h, adj);
    for (int c : adj[v]) {
        if (c != parent[v] && c != heavy[v])
            decompose(c, c, adj);
    }
}

// se asume v=0 raíz
void init(vector<vector<int>> const& adj) {
    int n = adj.size();
    parent = vector<int>(n);
    depth = vector<int>(n);
    heavy = vector<int>(n, -1);
    head = vector<int>(n);
    pos = vector<int>(n);
    cur_pos = 0;

    dfs(0, adj);
    decompose(0, 0, adj);
}

// max-query en este ejemplo
// Encontramos lca(a, b)=L, hacemos la query del segment tree a->L y L->b
// Como hay un segment tree por heavy path y hay logn heavy paths, O(log^2 n)
// |- (pero si precalculas los prefijos es logn)
// TODO: implementar segment_tree_query(-, -). En lugar de implementar logn
//       segment-trees, se puede implementar uno solo con segmentos disjuntos
//       por cada heavy path
int query(int a, int b) {
    int res = 0;
    for (; head[a] != head[b]; b = parent[head[b]]) {
        if (depth[head[a]] > depth[head[b]])
            swap(a, b);
        int cur_heavy_path_max = segment_tree_query(pos[head[b]], pos[b]);
        res = max(res, cur_heavy_path_max);
    }
    if (depth[a] > depth[b])
        swap(a, b);
    int last_heavy_path_max = segment_tree_query(pos[a], pos[b]);
    res = max(res, last_heavy_path_max);
    return res;
}

\end{minted}

	\needspace{1\baselineskip}
	{
		\phantomsection
		\subsubsection*{LCA}
		\addcontentsline{toc}{subsubsection}{LCA}
	}
	\begin{minted}{cpp}
// Puede ser útil recordar que dist(a,b) = depth(a) + depth(b) - 2*depth(lca(a,b))
int n, l;                // n = número de nodos, l = log2(n)
vector<vector<int>> adj; // Lista de adyacencia del árbol

int timer;              // Temporizador para asignar tin y tout
vector<int> tin, tout;  // Tiempo de entrada y salida en DFS
vector<vector<int>> up; // up[v][i] = 2^i-ésimo ancestro de v

// DFS para preprocesar ancestros y tiempos
void dfs(int v, int p)
{
  tin[v] = ++timer; // Al entrar al nodo v
  up[v][0] = p;     // 2^0-ésimo ancestro es el padre directo

  // Computar los 2^i-ésimos ancestros
  for (int i = 1; i <= l; ++i)
    up[v][i] = up[up[v][i - 1]][i - 1];

  // Continuar DFS con los hijos
  for (int u : adj[v])
  {
    if (u != p)
      dfs(u, v);
  }

  tout[v] = ++timer; // Al salir del nodo v
}

// Verifica si u es ancestro de v usando tin y tout
bool is_ancestor(int u, int v)
{
  return tin[u] <= tin[v] && tout[u] >= tout[v];
}

// Devuelve el Lowest Common Ancestor de u y v
int lca(int u, int v)
{
  // Si uno es ancestro del otro, es el LCA
  if (is_ancestor(u, v))
    return u;
  if (is_ancestor(v, u))
    return v;

  // Subimos a u hasta que su ancestro está justo por debajo del LCA
  for (int i = l; i >= 0; --i)
  {
    if (!is_ancestor(up[u][i], v))
      u = up[u][i];
  }

  return up[u][0]; // El padre de u en este punto es el LCA
}

// Preprocesamiento inicial para Binary Lifting
void preprocess(int root)
{
  tin.resize(n);
  tout.resize(n);
  timer = 0;
  l = ceil(log2(n)); // Límite de profundidad de los ancestros
  up.assign(n, vector<int>(l + 1));
  dfs(root, root); // Empezamos el DFS desde la raíz
}

int main()
{
  n = 7;
  adj.resize(n);
  adj[0] = {1, 2};
  adj[1] = {0, 3, 4};
  adj[2] = {0, 5, 6};
  adj[3] = {1};
  adj[4] = {1};
  adj[5] = {2};
  adj[6] = {2};

  preprocess(0); // Raíz del árbol en 0

  cout << lca(3, 4) << endl; // Salida esperada: 1
  cout << lca(5, 6) << endl; // Salida esperada: 2
  cout << lca(3, 5) << endl; // Salida esperada: 0
}

\end{minted}

	\needspace{5\baselineskip}
	{
		\phantomsection
		\section*{Strings}
		\markboth{STRINGS}{}
		\addcontentsline{toc}{section}{Strings}
	}
	{
		\phantomsection
		\subsection*{Aho-Corasik}
		\addcontentsline{toc}{subsection}{Aho-Corasik}
	}
	\begin{minted}{cpp}
const int M = 26;
struct node
{
  vector<int> child;
  int p = -1;
  char c = 0;
  int suffixLink = -1, endLink = -1;
  int id = -1;
  node(int p = -1, char c = 0) : p(p), c(c)
  {
    child.resize(M, -1);
  }
};

struct AhoCorasick
{
  vector<node> t;
  vector<int> lengths; // Longitud de las palabras en el trie
  vector<int> occ;     // Ocurrencias por patrón
  int wordCount = 0;
  AhoCorasick()
  {
    t.emplace_back();
  }

  // Regresa el ID que se le asigno a la palabra agregada (0-indexado)
  int add(const string &s)
  {
    int u = 0; // Raíz temporal
    for (char c : s)
    {
      if (t[u].child[c - 'a'] == -1)
      {
        t[u].child[c - 'a'] = t.size();
        t.emplace_back(u, c);
      }
      u = t[u].child[c - 'a'];
    }
    if (t[u].id == -1) // Por si se trata de agregra una palabra que ya estaba
      t[u].id = wordCount++;
    lengths.push_back(s.size());
    return t[u].id;
  }

  void link(int u)
  {
    if (u == 0) // Si eres la raíz
    {
      t[u].suffixLink = 0;
      t[u].endLink = 0;
      return;
    }
    if (t[u].p == 0) // Si tu padre es la raíz
    {
      t[u].suffixLink = 0;
      if (t[u].id != -1)
        t[u].endLink = u;
      else
        t[u].endLink = t[t[u].suffixLink].endLink;
      return;
    }

    int v = t[t[u].p].suffixLink;
    char c = t[u].c;
    while (true)
    {
      if (t[v].child[c - 'a'] != -1)
      {
        t[u].suffixLink = t[v].child[c - 'a'];
        break;
      }
      if (v == 0)
      {
        t[u].suffixLink = 0;
        break;
      }
      v = t[v].suffixLink;
    }
    if (t[u].id != -1) // Si eres una palabra
      t[u].endLink = u;
    else
      t[u].endLink = t[t[u].suffixLink].endLink;
  }

  // BFS para construir los suffix links y end lins. Llamar a esto
  // cuando ya hayas metido todas strings del diccionario al trie
  void build()
  {
    queue<int> Q;
    Q.push(0);
    while (!Q.empty())
    {
      int u = Q.front();
      Q.pop();
      link(u);
      for (int v = 0; v < M; ++v)
        if (t[u].child[v] != -1)
          Q.push(t[u].child[v]);
    }
    occ.assign(wordCount, 0);
  }

  int match(const string &text)
  {
    int u = 0;
    int ans = 0;
    for (int j = 0; j < text.size(); ++j)
    {
      int i = text[j] - 'a';
      while (true)
      {
        if (t[u].child[i] != -1)
        {
          u = t[u].child[i];
          break;
        }
        if (u == 0)
          break;
        u = t[u].suffixLink;
      }

      // Temporal de u para moverse por el endLink para el caso
      // en el que haya sufijos propios de lo que llevamos visitado
      // que su vez sean palabras
      int v = u;
      while (true)
      {
        v = t[v].endLink;
        if (v == 0)
          break;
        ++ans;
        // int idx = j + 1 - lengths[t[v].id];
        // cout << "Found word #" << t[v].id << " at position " << idx << "\n";
        occ[t[v].id]++;
        v = t[v].suffixLink; // Para eviar el caso en el que el nodo de v es el de una palabra
      }
    }
    return ans;
  }
};

int main()
{
  AhoCorasick ac;
  string text;
  cin >> text;
  int n;
  cin >> n;
  vector<int> ids(n); // Por si se agrega varias veces la misma palabra
  for (int i = 0; i < n; i++)
  {
    string word;
    cin >> word;
    ids[i] = ac.add(word);
  }
  ac.build();
  ac.match(text);
  for (int i = 0; i < n; i++)
    cout << ac.occ[ids[i]] << endl;
  return 0;
}

\end{minted}

	\needspace{2\baselineskip}
	{
		\phantomsection
		\subsection*{Counting Substrings}
		\addcontentsline{toc}{subsection}{Counting Substrings}
	}
	\begin{minted}{cpp}

/*
Given a string s and q queries, each query is a string t
For each request you need to count how many times the string t
occurs as a substring in s.
O(t log(n)) where n = s.size() for each query 
*/

int lwb(vi &SA, string &s, string &t)
{
  int l = 0;
  int r = SA.size() - 1;

  int res = SA.size();
  while (l <= r)
  {
    int m = l + (r - l) / 2;
    if (s.substr(SA[m], t.size()) >= t)
    {
      res = m;
      r = m - 1;
    }
    else
    {
      l = m + 1;
    }
  }
  return res;
}

int upb(vi &SA, string &s, string &t)
{
  int l = 0;
  int r = SA.size() - 1;
  int res = SA.size();
  while (l <= r)
  {
    int m = l + (r - l) / 2;
    if (s.substr(SA[m], t.size()) > t)
    {
      res = m;
      r = m - 1;
    }
    else
    {
      l = m + 1;
    }
  }
  return res;
}

int main()
{
  ios::sync_with_stdio(0);
  cin.tie(nullptr);

  string s;
  cin >> s;
  s += "$";

  vi SA = Suffix_Array(s);

  int q;
  cin >> q;

  while (q--)
  {
    string t;
    cin >> t;
    cout << upb(SA, s, t) - lwb(SA, s, t) << endl;
  }
  return 0;
}
\end{minted}

	\needspace{2\baselineskip}
	{
		\phantomsection
		\subsection*{Finding Borders}
		\addcontentsline{toc}{subsection}{Finding Borders}
	}
	\begin{minted}{cpp}
/*
A border of a string is a prefix that is also a suffix of the string but not the whole string. For example, the borders of abcababcab are ab and abcab.
Your task is to find all border lengths of a given string: 
*/
void Finding_Borders(string& s){
  vector<int> pi = prefix_function(s);
  // Empezamos con el borde más grande
  int j = pi[s.size() - 1];
  vector<int> res;
  while (j > 0)
  {
    res.push_back(j);
    // El siguiente borde más grande 
    j = pi[j - 1];
  }
  for (int i = res.size() - 1; i >= 0; i--)
    cout << res[i] << " ";
}
\end{minted}

	\needspace{2\baselineskip}
	{
		\phantomsection
		\subsection*{Finding Periods}
		\addcontentsline{toc}{subsection}{Finding Periods}
	}
	\begin{minted}{cpp}
/*
A period of a string is a prefix that can be used to generate the whole string by repeating the prefix. The last repetition may be partial. For example, the periods of abcabca are abc, abcabc and abcabca.
Find all period lengths of a string: 
*/
void Finding_Periods(string& s){
  int n = s.size();
  vector<int> z = Z_function(s);

  for (int i = 0; i < n; i++)
    if (z[i] == n - i)
      cout << i << " ";

  cout << n << endl;
}
\end{minted}

	\needspace{2\baselineskip}
	{
		\phantomsection
		\subsection*{Hashing}
		\addcontentsline{toc}{subsection}{Hashing}
	}
	\begin{minted}{cpp}
struct Hash
{
  ll P = 1777771; // Primo para sacar el hash
  ll MOD[2], PI[2];
  vector<ll> h[2]; // Contiene los prefix sum de los hahses  
  vector<ll> pi[2]; // pi[k][i] contiene PI[k] elevado a la i 
  Hash(string &s)
  {
    MOD[0] = 999727999;
    MOD[1] = 1070777777;
    PI[0] = 325255434; // Inverso multiplicativo de P modulo MOD[0]
    PI[1] = 10018302; // Inverso multiplicativo de P modulo MOD[1]
    for (int k = 0; k < 2; k++)
      h[k].resize(s.size() + 1), pi[k].resize(s.size() + 1);
    for (int k = 0; k < 2; k++)
    {
      h[k][0] = 0;
      pi[k][0] = 1;
      ll p = 1;
      for (int i = 1; i < s.size() + 1; i++)
      {
        h[k][i] = (h[k][i - 1] + p * s[i - 1]) % MOD[k];
        pi[k][i] = (1LL * pi[k][i - 1] * PI[k]) % MOD[k];
        p = (p * P) % MOD[k];
      }
    }
  }
  // Returns the hash of substring [s,e)
  ll get(int s, int e)
  {
    ll h0 = (h[0][e] - h[0][s] + MOD[0]) % MOD[0];
    h0 = (1LL * h0 * pi[0][s]) % MOD[0];
    ll h1 = (h[1][e] - h[1][s] + MOD[1]) % MOD[1];
    h1 = (1LL * h1 * pi[1][s]) % MOD[1];
    return (h0 << 32) | h1;
  }
};

// O(n²)
int Distinct_Substrings(string& s){
  int n = s.size();
  Hash s_hash(s);
  unordered_set<ll> aux;
  for (int i = 0; i < n; i++)
    for (int j = i; j < n; j++)
      aux.insert(s_hash.get(i, j + 1));

 return aux.size(); 
}

// Nos dice si la substring s[l2...r2] es menor que la string s[l1...r1] en
// tiempo log_2(n)
bool less_str(string &s, Hash &hash, int l1, int r1, int l2, int r2)
{
  int l = 0;
  int r = s.size() - 1;

  // Busqueda binaria en la longitud máxima del prefijo que es igual en ambas strings
  int last = 0;
  while (l <= r)
  {
    int m = l + (r - l) / 2;
    int h1 = hash.get(l1, l1 + m);
    int h2 = hash.get(l2, l2 + m);

    if (h1 == h2)
    {
      last = m;
      l = m + 1;
    }
    else
    {
      r = m - 1;
    }
  }

  // s[l1...l1+last) == s[l2...l2+last)
  // Si last < s.size() && s[l2+last] < s[l1+last] entonces
  // s[l2...r2] es menor que la string s[l1...r1]
  return last < s.size() && s[l2 + last] < s[l1 + last];
}
\end{minted}

	\needspace{2\baselineskip}
	{
		\phantomsection
		\subsection*{Hashing Strings as vectors}
		\addcontentsline{toc}{subsection}{Hashing Strings as vectors}
	}
	\begin{minted}{cpp}
vector<ll> P = {325255434, 10018302};
vector<ll> M = {999727999, 1070777777};

/*
  Para tratar todas las permutaciones de una misma cadena como equivalentes,
  se puede representar la cadena mediante su vector de frecuencias.
  Por ejemplo, la cadena "abacb" se representa como:
    f = {2, 2, 1, 0, 0, ...}
         a  b  c  d  e ...
  donde f[i] indica cuántas veces aparece el carácter en la posición i del alfabeto en la string dada.

  Para comparar estas representaciones de forma eficiente, se puede calcular
  un hash del vector de frecuencias. Sea:
    - f el vector de frecuencias de una cadena,
    - |A| el tamaño del alfabeto (por ejemplo, 26 para letras minúsculas),
    - b y M constantes enteras elegidas para el cálculo del hash,

  entonces definimos el hash como:
    hash(f) = (f[0] * b^0 + f[1] * b^1 + ... + f[|A|-1] * b^|A|-1) mod M

  Nota: Si ya se ha calculado el hash del vector correspondiente a la subcadena s[L...R] entonces calcular el hash del vector de la subcadena definida por 
  s[L...R+1] implica:
    hash(S[L...R+1]) = ( hash(S[L...R]) + b^(pos(S[R+1])) ) mod M 
  donde pos(S[R+1]) es la posición del caracter S[R+1] en el alfabeto. 
*/
void Hashing_Strigns_as_Vectors()
{
  string s; 
  int sigma = 27; // Tamaño del alfabeto
  int n = s.size();

  vector<vector<ll>> pot(2, vector<ll>(sigma, 0));
  pot[0][0] = pot[1][0] = 1;

  for (int i = 1; i < sigma; i++)
  {
    pot[0][i] = pot[0][i - 1] * P[0] % M[0];
    pot[1][i] = pot[1][i - 1] * P[1] % M[1];
  }

  // Calcula el hash de todos los vectores de frecuencias de 
  // todas las substrings de s
  for (int i = 0; i < n; i++)
  {
    ll hash1 = 0;
    ll hash2 = 0;
    for (int j = i; j < n; j++)
    {
      int pos_c = s[j] - 'a';
      hash1 += pot[0][pos_c] % M[0];
      hash2 += pot[1][pos_c] % M[1];

      ll final_hash = (hash1 << 32) | hash2;
    }
  }
}


\end{minted}

	\needspace{2\baselineskip}
	{
		\phantomsection
		\subsection*{KMP}
		\addcontentsline{toc}{subsection}{KMP}
	}
	\begin{minted}{cpp}
vector<int> KMP(const string &text, const string &pat)
{
  int n = text.size();
  vector<int> pi = prefix_function(pat);
  int j = pi[0];
  vector<int> pi_text(n, 0);
  for (int i = 0; i < n; i++)
  {
    while (j > 0 && text[i] != pat[j]) j = pi[j - 1];
    if (text[i] == pat[j]) j++;
    pi_text[i] = j;
    if (j == pat.size()) j = pi[j - 1];
  }
  // Las ocurrencias de pat.size() en pi_text
  // son las apariciones de pat en text
  return pi_text;
}
\end{minted}

	\needspace{2\baselineskip}
	{
		\phantomsection
		\subsection*{Longest Common Substring}
		\addcontentsline{toc}{subsection}{Longest Common Substring}
	}
	\begin{minted}{cpp}

// Nos dice si i pertenece a la string a y j a la string b o viceversa
bool distintos(int n1, int n, int i, int j)
{
  return (i >= 0 && i < n1 && j >= n1 + 1 && j < n) || (j >= 0 && j < n1 && i >= n1 + 1 && i < n);
}

string LCS()
{
  string a;
  cin >> a;
  int n1 = a.size();

  string b;
  cin >> b;
  int n2 = b.size();

  string s = a + "$" + b + "#";

  pair<vi, vi> result = Suffix_Array_and_LCP(s);
  vi SA = result.first;
  vi LCP = result.second;

  int res = 0;
  int maximo = INT_MIN;
  for (int i = 1; i < s.size(); i++)
  {
    if (distintos(n1, s.size(), SA[i], SA[i - 1]))
    {
      if (LCP[SA[i]] > maximo)
      {
        maximo = LCP[SA[i]];
        res = SA[i];
      }
    }
  }
  return s.substr(res, LCP[res]); 
}
\end{minted}

	\needspace{2\baselineskip}
	{
		\phantomsection
		\subsection*{Manacher}
		\addcontentsline{toc}{subsection}{Manacher}
	}
	\begin{minted}{cpp}
// Palíndromos impares centrados en cada posición
vector<int> manacher_odd(const string& s) {
    int n = s.size();
    vector<int> odd(n);
    int l = 0, r = -1;
    for (int i = 0; i < n; ++i) {
        int k = (i > r ? 1 : min(odd[l + r - i], r - i + 1));
        while (i - k >= 0 && i + k < n && s[i - k] == s[i + k]) k++;
        odd[i] = k--;
        if (i + k > r) {
            l = i - k;
            r = i + k;
        }
    }
    return odd;
}
// Palíndromos pares centrados entre posiciones
vector<int> manacher_even(const string& s) {
    int n = s.size();
    vector<int> even(n);
    int l = 0, r = -1;
    for (int i = 0; i < n; ++i) {
        int k = (i > r ? 0 : min(even[l + r - i + 1], r - i + 1));
        while (i - k - 1 >= 0 && i + k < n && s[i - k - 1] == s[i + k]) k++;
        even[i] = k--;
        if (i + k > r) {
            l = i - k - 1;
            r = i + k;
        }
    }
    return even;
}

// Función para contar todos los palíndromos
int count_palindromes(const string& s) {
    vector<int> odd = manacher_odd(s);
    vector<int> even = manacher_even(s);
    int total_palindromes = 0;
    // Contar los palíndromos impares
    for (int i = 0; i < s.size(); ++i) {
        total_palindromes += odd[i];  // Cada odd[i] indica cuántos palíndromos impares centrados en i
    }
    // Contar los palíndromos pares
    for (int i = 0; i < s.size(); ++i) {
        total_palindromes += even[i]; // Cada even[i] indica cuántos palíndromos pares centrados entre i-1 e i
    }
    return total_palindromes;
}

int main() {
    string s = "abacaba";
    int total = count_palindromes(s);
    cout << "Total de palíndromos en la cadena: " << total << endl;
    return 0;
}

\end{minted}

	\needspace{2\baselineskip}
	{
		\phantomsection
		\subsection*{Number of Balanced Substrings}
		\addcontentsline{toc}{subsection}{Number of Balanced Substrings}
	}
	\begin{minted}{cpp}
// Cuenta el número de subcadenas que están en 1-Dyck en O(n)
// (subcadenas que están en el leng. de paréntesis balanceados)

ll aux_subdyck(const vector<int>& to_open, int start, int end) // auxiliar
{
    ll streak = 0, acc = 0; // acc es la cuenta, streak es la racha de 
                             // paréntesis adyacentes en el mismo nivel
    for (int i = end-1; i >= start; --i)
    {
        int op = to_open[i];
        if (op < 0) { acc += (streak*(streak+1))>>1; streak = 0; }
        else { streak++; acc += aux_subdyck(to_open, op+1, i); i = op; }
    }
    acc += (streak*(streak+1))>>1;
    return acc;
}

ll cnt_subdyck(const string& s, char opensymb, char closesymb) // O(n) dp
{
    int n = s.length();
    vector<int> opens; opens.reserve(n); // pila de idx de parens abiertos
    vector<int> to_open(n,-1); // dónde se encuentra la pareja '(' de cada ')'
    for (int i = 0; i < n; i++)
    {
        if (s[i] == opensymb)
            opens.push_back(i);
        else if (s[i] == closesymb and opens.size())
        {
            to_open[i] = opens[opens.size()-1];
            opens.pop_back();
        }
    }
    return aux_subdyck(to_open, 0, n); // cuenta por niveles de parens
}

\end{minted}

	\needspace{2\baselineskip}
	{
		\phantomsection
		\subsection*{Number of Distinct Substrings}
		\addcontentsline{toc}{subsection}{Number of Distinct Substrings}
	}
	\begin{minted}{cpp}
ll Distinct_Substrings()
{
  string s;
  cin >> s;
  s += "$";

  pair<vi, vi> result = Suffix_Array_and_LCP(s);
  vi SA = result.first;
  vi LCP = result.second;

  ll sum = 0;
  ll n = s.size() - 1; // Para no contar al '$' agregado

  for (int i = 0; i < n; i++)
    sum += LCP[i];

  ll number_of_substrings = (n * (n + 1)) / 2;
  return number_of_substrings - sum; 
}
\end{minted}

	\needspace{2\baselineskip}
	{
		\phantomsection
		\subsection*{Occurrences Of Each Prefixes}
		\addcontentsline{toc}{subsection}{Occurrences Of Each Prefixes}
	}
	\begin{minted}{cpp}
vector<int> occurrences_of_each_prefixes(const string& s){
  int n = s.size();
  vector<int> pref = prefix_function(s);
  vector<int> occ(n + 1);
  for (int i = 0; i < n; i++)
    occ[pref[i]]++;
  for (int i = n - 1; i > 0; i--)
    occ[pref[i - 1]] += occ[i];
  for (int i = 0; i <= n; i++)
    occ[i]++;
  return occ; 
}
\end{minted}

	\needspace{2\baselineskip}
	{
		\phantomsection
		\subsection*{Prefix Function}
		\addcontentsline{toc}{subsection}{Prefix Function}
	}
	\begin{minted}{cpp}
/*
pi[i] = Longitud del prefijo propio más grande
en la substring s[0...i] que también es sufijo de 
esa misma substring. Por definición pi[0] = 0. 

s =  b a b a b a c b b a 
pi = 0 0 1 2 3 4 0 1 1 2 
    ----------------------
     0 1 2 3 4 5 6 7 8 9 
*/
vector<int> prefix_function(const string &s)
{
  int n = s.size();
  vector<int> pi(n, 0);
  int j = 0;
  for (int i = 1; i < n; i++)
  {
    while (j > 0 && s[i] != s[j]) j = pi[j - 1];
    if (s[i] == s[j]) j++;
    pi[i] = j;
  }
  return pi;
}

/*
A border of a string is a prefix that is also a suffix of the string but not the whole string. For example, the borders of abcababcab are ab and abcab.
Your task is to find all border lengths of a given string: 
*/
void Finding_Borders(string& s){
  vector<int> pi = prefix_fun(s);
  // Empezamos con el borde más grande
  int j = pi[s.size() - 1];
  vector<int> res;
  while (j > 0)
  {
    res.push_back(j);
    // El siguiente borde más grande 
    j = pi[j - 1];
  }
  for (int i = res.size() - 1; i >= 0; i--)
    cout << res[i] << " ";
}
\end{minted}

	\needspace{2\baselineskip}
	{
		\phantomsection
		\subsection*{Repeating Substring}
		\addcontentsline{toc}{subsection}{Repeating Substring}
	}
	\begin{minted}{cpp}
/*
A repeating substring is a substring that occurs in two (or more) locations in the string. Your task is to find the longest repeating substring in a given string.
*/
string Repeating_Substring()
{
  string s;
  cin >> s;
  s += "$";
  pair<vi, vi> result = Suffix_Array_and_LCP(s);
  vi SA = result.first;
  vi LCP = result.second;

  int ans = 0;
  for (int i = 0; i < s.size(); i++)
    ans = max(ans, LCP[SA[i]]);

  if (ans == 0)
    return ""; 

  for (int i = 0; i < s.size(); i++)
    if (ans == LCP[SA[i]])
      return s.substr(SA[i], LCP[SA[i]]);
    
  return "";
}
\end{minted}

	\needspace{2\baselineskip}
	{
		\phantomsection
		\subsection*{Suffix Array and LCP}
		\addcontentsline{toc}{subsection}{Suffix Array and LCP}
	}
	\begin{minted}{cpp}
/*
Suffix array (SA):
Let s be a string of length n.The i-th suffix of s is the substring
s[i...n - 1].
A suffix array will contain integers that represent the starting indexes of the all the suffixes of a given string, after the aforementioned suffixes are sorted.

Note:  The order of the sorted cyclic shifts is equivalent to the order of the sorted suffixes

LCP (Longest Common Prefix) array:
A value lcp[SA[i]] indicates length of the longest common prefix of the suffixes indexed in SA[i] and SA[i+1]. 

txt[0..n-1] = "b a n a n a $"
               0 1 2 3 4 5 6

suffix[]  = {6,    5,   3,     1,       0,        4,    2}
          = {"$", "a", "ana", "anana", "banana", "na", "nana"}
            ---------------------------------------------------
             0     1    2      3        4         5     6

for (int i = 1; i < n; i++)
    cout << LCP[SA[i]] << " "; 
lcp        = {0, 1, 3, 0, 0, 2, 0}
lcp[SA[0]] = LCP of "$" and "a"          = 0
lcp[SA[1]] = LCP of "a" and "ana"        = 1
lcp[SA[2]] = LCP of "ana" and "anana"    = 3
lcp[SA[3]] = LCP of "anana" and "banana" = 0
lcp[SA[4]] = LCP of "banana" and "na"    = 0
lcp[SA[5]] = LCP of "na" and "nana"      = 2
*/

void countingSort(vi &SA, vi &eq)
{
  vi count(sz(eq));

  for (int i = 0; i < sz(eq); i++)
    count[eq[i]]++;

  for (int i = 1; i < sz(eq); i++)
    count[i] += count[i - 1];

  vi out(sz(eq));

  for (int i = sz(eq) - 1; i >= 0; i--)
  {
    int val = eq[SA[i]];
    out[count[val] - 1] = SA[i];
    count[val]--;
  }
  SA = out;
}

pair<vi, vi> Suffix_Array_and_LCP(const string &s)
{
  int n = sz(s);
  vector<pair<char, int>> aux(n);

  for (int i = 0; i < n; i++)
    aux[i] = {s[i], i};

  sort(all(aux));

  vi SA(n);

  for (int i = 0; i < n; i++)
    SA[i] = aux[i].second;

  vi eq(n);

  eq[SA[0]] = 0;
  for (int i = 1; i < n; i++)
    eq[SA[i]] = (aux[i - 1].first != aux[i].first) ? eq[SA[i - 1]] + 1 : eq[SA[i - 1]];

  for (int h = 0; (1 << h) < n; h++)
  {
    int k = 1 << h;

    for (int i = 0; i < n; i++)
      SA[i] = (SA[i] - k + n) % n;

    countingSort(SA, eq);

    vi aux(n);
    aux[SA[0]] = 0;
    for (int i = 1; i < n; i++)
    {
      pair<int, int> ant = {eq[SA[i - 1]],
                            eq[(SA[i - 1] + k) % n]};
      pair<int, int> act = {eq[SA[i]], eq[(SA[i] + k) % n]};
      aux[SA[i]] = (ant != act) ? aux[SA[i - 1]] + 1 : aux[SA[i - 1]];
    }

    eq = aux;
  }

  vi LCP(n);

  int k = 0;
  for (int i = 0; i < n - 1; i++)
  {
    int j = SA[eq[i] - 1];

    while (s[i + k] == s[j + k])
      k++;

    LCP[i] = k;

    k = max(k - 1, 0);
  }

  return {SA, LCP};
}

int main()
{
  ios::sync_with_stdio(0);
  cin.tie(nullptr);
  string s;
  cin >> s;
  s += "$";
  pair<vi, vi> result = Suffix_Array_and_LCP(s);
  vi SA = result.first;
  vi LCP = result.second;
  return 0;
}
\end{minted}

	\needspace{2\baselineskip}
	{
		\phantomsection
		\subsection*{Trie}
		\addcontentsline{toc}{subsection}{Trie}
	}
	\begin{minted}{cpp}
struct Node
{
  bool isWord = false;
  map<char, Node *> letters;
};
struct Trie
{
  Node *root;
  Trie()
  {
    root = new Node();
  }
  inline bool exists(Node *actual, const char &c)
  {
    return actual->letters.find(c) != actual->letters.end();
  }
  void InsertWord(const string &word)
  {
    Node *current = root;
    for (auto &c : word)
    {
      if (!exists(current, c))
        current->letters[c] = new Node();
      current = current->letters[c];
    }
    current->isWord = true;
  }
  bool FindWord(const string &word)
  {
    Node *current = root;
    for (auto &c : word)
    {
      if (!exists(current, c))
        return false;
      current = current->letters[c];
    }
    return current->isWord;
  }
  void printRec(Node *actual, string acum)
  {
    if (actual->isWord)
    {
      cout << acum << "\n";
    }
    for (auto &next : actual->letters)
      printRec(next.second, acum + next.first);
  }
  void printWords(const string &prefix)
  {
    Node *actual = root;
    for (auto &c : prefix)
    {
      if (!exists(actual, c))
        return;
      actual = actual->letters[c];
    }
    printRec(actual, prefix);
  }
};
\end{minted}

	\needspace{2\baselineskip}
	{
		\phantomsection
		\subsection*{Z function}
		\addcontentsline{toc}{subsection}{Z function}
	}
	\begin{minted}{cpp}
/**
z[i] = longitud de la substring más grande que empieza en i 
y que también es un prefijo de la string s 

s = a a b x a a y a a b 
z = 0 1 0 0 2 1 0 3 1 0
  -----------------------
    0 1 2 3 4 5 6 7 8 9 
*/
vector<int> Z_function(const string &s)
{
  int n = s.size();
  vector<int> z(n, 0);
  int l = 0, r = 0; 
  z[0] = 0;
  for (int i = 1; i < n; i++)
  {
    if (i < r) z[i] = min(r - i, z[i - l]);
    while (i + z[i] < n && s[z[i]] == s[i + z[i]]) z[i]++;
    if (i + z[i] > r)
    {
      l = i;
      r = i + z[i];
    }
  }
  return z;
}


\end{minted}

	\needspace{5\baselineskip}
	{
		\phantomsection
		\section*{Dynamic Programming}
		\markboth{DINAMIC PROGRAMMING}{}
		\addcontentsline{toc}{section}{Dynamic Programming}
	}
	{
		\phantomsection
		\subsection*{Coin Combinations I}
		\addcontentsline{toc}{subsection}{Coin Combinations I}
	}
	\begin{minted}{cpp}
const int X = 1e6 + 1;
ll dp[X];
ll coin_combinations_I()
{
  int n, x; cin >> n >> x;
  if (n <= 0 || x < 0) return 0;

  vector<ll> coins(n);
  for (int i = 0; i < n; i++)
      cin >> coins[i];
  dp[0] = 1;

  for (int sum = 1; sum <= x; sum++)
    for (int c : coins)
      if (sum - c >= 0)
        {
          dp[sum] += dp[sum - c];
          dp[sum] %= MOD;
        }
  return dp[x];
}


\end{minted}

	\needspace{2\baselineskip}
	{
		\phantomsection
		\subsection*{Coin Combinations II}
		\addcontentsline{toc}{subsection}{Coin Combinations II}
	}
	\begin{minted}{cpp}
const int X = 1e6 + 1;
ll dp[X];
ll coin_combinations_II()
{
  int n, x; cin >> n >> x;
  if (n <= 0 || x < 0) return 0;

  vector<ll> coins(n);
  for (int i = 0; i < n; i++)
    cin >> coins[i];
  dp[0] = 1;

  for (int c : coins)
    for (int sum = 0; sum <= x; sum++)
      if (sum - c >= 0)
      {
        dp[sum] += dp[sum - c];
        dp[sum] %= MOD;
      }
  return dp[x];
}

\end{minted}

	\needspace{2\baselineskip}
	{
		\phantomsection
		\subsection*{Digit DP}
		\addcontentsline{toc}{subsection}{Digit DP}
	}
	\begin{minted}{cpp}
#include <bits/stdc++.h>
using namespace std;

using ll = long long;

const int N = 19;
const int DIG = 10;
const int CEROS = 2;
const int LIBRE = 2;

ll dp[N][DIG][CEROS][LIBRE];

ll counting(string &numero, int n, int ant, bool ceros, bool libre)
{
  if (n == 0)
    return 1;

  int dig = numero[numero.size() - n] - '0';

  if (ant != -1 && dp[n][ant][ceros][libre] != -1)
    return dp[n][ant][ceros][libre];

  ll res = 0;
  int tope = (libre) ? 9 : dig;

  for (int j = 0; j <= tope; j++)
  {

    if (j == ant && !ceros)
      continue;

    res += counting(numero, n - 1, j, ceros && j == 0, libre || j < dig);
  }

  if (ant != -1)
    dp[n][ant][ceros][libre] = res;
  return res;
}

// Your task is to count the number of integers between a and b where no two adjacent digits are the same.
int main()
{
  ios::sync_with_stdio(0);
  cin.tie(nullptr);

  ll a, b;
  cin >> a >> b;
  a--;
  memset(dp, -1, sizeof(dp));
  string b_s = to_string(b);
  ll res = counting(b_s, b_s.size(), -1, 1, 0);

  if (a >= 0)
  {
    memset(dp, -1, sizeof(dp));
    string a_s = to_string(a);
    res -= counting(a_s, a_s.size(), -1, 1, 0);
  }

  cout << res << endl;
  return 0;
}
\end{minted}

	\needspace{2\baselineskip}
	{
		\phantomsection
		\subsection*{Edit Distance}
		\addcontentsline{toc}{subsection}{Edit Distance}
	}
	\begin{minted}{cpp}

const int N = 5000+1; 
const int M = 5000+1; 
int dp[N][M]; 
int edit_distance()
{
  string s, t; cin >> s >> t;
  int n = s.size(); int m = t.size();

  for (int j = m; j >= 0; j--)
    dp[n][j] = t.size() - j;
  for (int i = n; i >= 0; i--)
    dp[i][m] = s.size() - i;
  
  for (int i = n - 1; i >= 0; i--)
    for (int j = m - 1; j >= 0; j--)
      if (s[i] == t[j])
        dp[i][j] = dp[i + 1][j + 1];
      else
        dp[i][j] = min({1 + dp[i + 1][j + 1], 
                        1 + dp[i + 1][j], 
                        1 + dp[i][j + 1]});
    
  return dp[0][0]; 
}

\end{minted}

	\needspace{2\baselineskip}
	{
		\phantomsection
		\subsection*{Knapsack 01}
		\addcontentsline{toc}{subsection}{Knapsack 01}
	}
	\begin{minted}{cpp}
const ll W = 1e5 + 1;
const ll N = 100 + 1;
ll value[N], weight[N], dp[N][W];
ll knapsack_01()
{
  int n, w; cin >> n >> w;
  for (int i = 0; i < n; i++)
  { cin >> weight[i]; cin >> value[i]; }

  for (int i = 1; i <= n; i++)
    for (int j = 0; j <= w; j++)
      if (j - weight[i - 1] < 0) 
        dp[i][j] = dp[i - 1][j];
      else 
        dp[i][j] = max(dp[i - 1][j], 
        dp[i - 1][j - weight[i - 1]] + value[i - 1]);
    
  return dp[n][w];
}

\end{minted}

	\needspace{2\baselineskip}
	{
		\phantomsection
		\subsection*{Knapsack 01 optimized}
		\addcontentsline{toc}{subsection}{Knapsack 01 optimized}
	}
	\begin{minted}{cpp}
const ll W = 1e5 + 1;
const ll N = 100 + 1;
ll value[N], weight[N], dp[W];
ll knapsack_01_op()
{
  ll n, w; cin >> n >> w;
  for (int i = 0; i < n; i++)
  { cin >> weight[i]; cin >> value[i]; }

  for (int i = 0; i < n; i++)
    for (int j = w; j >= 0; j--)
      if (j - weight[i] >= 0)
        dp[j] = max(dp[j], dp[j - weight[i]] + value[i]);
      
  return dp[w];
}

\end{minted}

	\needspace{2\baselineskip}
	{
		\phantomsection
		\subsection*{Knapsack 2}
		\addcontentsline{toc}{subsection}{Knapsack 2}
	}
	\begin{minted}{cpp}
const int N = 100 + 1;
const int MAX_V = N * 1e3 + 1;
ll dp[N][MAX_V], value[N], weight[N];
// Es como un coin change 
ll knpasack_01_2()
{
  int n, w; cin >> n >> w;
  for (int i = 0; i < n; i++)
  { cin >> weight[i]; cin >> value[i]; }

  for (int i = 0; i < N; i++)
    for (int j = 0; j < MAX_V; j++)
      dp[i][j] = INT_MAX;  
  
  dp[0][0] = 0;
  ll res = 0;
  for (int i = 1; i <= n; i++){
    for (int j = 0; j < MAX_V; j++)
    {
      if (j - value[i - 1] < 0)
        dp[i][j] = dp[i - 1][j];
      else
        dp[i][j] = min(dp[i - 1][j], 
                       dp[i - 1][j - value[i - 1]] + weight[i - 1]);
      if (dp[i][j] <= w)
        res = j;
    }
  }
  return res;
}


\end{minted}

	\needspace{2\baselineskip}
	{
		\phantomsection
		\subsection*{Knapsack 2 optimized}
		\addcontentsline{toc}{subsection}{Knapsack 2 optimized}
	}
	\begin{minted}{cpp}
const int N = 100 + 1;
const int MAX_V = N * 1e3 + 1;
ll dp[MAX_V], value[N], weight[N];
ll knpasack_01_2_op()
{
  int n, w; cin >> n >> w;

  for (int i = 0; i < n; i++)
  { cin >> weight[i]; cin >> value[i]; }

  for (int i = 0; i < MAX_V; i++)
    dp[i] = INT_MAX;
  
  dp[0] = 0;
  ll res = 0;
  for (int i = 1; i <= n; i++)
  {
    for (int j = MAX_V - 1; j >= 0; j--)
    {
      if (j - value[i - 1] >= 0)
        dp[j] = min(dp[j], dp[j - value[i - 1]] + weight[i - 1]);
      if (dp[j] <= w)
        res = max((ll)j, res);
    }
  }
  return res;
}

\end{minted}

	\needspace{2\baselineskip}
	{
		\phantomsection
		\subsection*{Longest Common Subsecuence}
		\addcontentsline{toc}{subsection}{Longest Common Subsecuence}
	}
	\begin{minted}{cpp}
int lcs(string &a, string &b)
{
  int m = a.size(), n = b.size();
  vector<vector<int>> aux(m + 1, vector<int>(n + 1));
  for (int i = 1; i <= m; ++i)
  {
    for (int j = 1; j <= n; ++j)
    {
      if (a[i - 1] == b[j - 1])
        aux[i][j] = 1 + aux[i - 1][j - 1];
      else
        aux[i][j] = max(aux[i - 1][j], aux[i][j - 1]);
    }
  }
  return aux[m][n];
}
\end{minted}

	\needspace{2\baselineskip}
	{
		\phantomsection
		\subsection*{Longest Increasing Subsecuence}
		\addcontentsline{toc}{subsection}{Longest Increasing Subsecuence}
	}
	\begin{minted}{cpp}
int LIS(vector<int> &nums)
{
  int n = nums.size();
  vector<int> lis_aux;
  for (int i = 0; i < n; i++)
  {
    int x = nums[i];
    int pos = lower_bound(lis_aux.begin(), lis_aux.end(), x) - lis_aux.begin();
    if (pos == lis_aux.size())
      lis_aux.push_back(x);
    else
      lis_aux[pos] = x;
  }
  return lis_aux.size(); 
}
\end{minted}

	\needspace{5\baselineskip}
	{
		\phantomsection
		\section*{Geometry}
		\markboth{GEOMETRY}{}
		\addcontentsline{toc}{section}{Geometry}
	}
	{
		\phantomsection
		\subsection*{Closest Pair}
		\addcontentsline{toc}{subsection}{Closest Pair}
	}
	\begin{minted}{cpp}
// encuentra la pareja de puntos más cercana en O(nlogn)

#include <bits/stdc++.h>
using namespace std;
#define ALL(x) x.begin(), x.end()

// se pueden cambiar para que funcione con double o float
using fl = long double;
#define HYPOT hypotl
#define ABS fabsl

struct pt;
using vpt = vector<pt>;
using ppt = pair<pt,pt>;
constexpr fl INF = numeric_limits<fl>::infinity();

struct pt { fl x, y; }; // aquí se podría poner más info extra a según
bool cmpx(const pt& a, const pt& b) { return a.x < b.x; }
bool cmpy(const pt& a, const pt& b) { return a.y < b.y; }
fl dist(const pt& a, const pt& b) { return HYPOT(a.x-b.x, a.y-b.y); }

// se usa para escoger pares
inline const pair<ppt,fl>& best_pair(const pair<ppt,fl>& a, const pair<ppt,fl>& b)
{
    return b.second < a.second ? b : a;
}

// para crear pares de puntos con su distancia
inline pair<ppt, fl> mk_p(const pt& a, const pt& b)
{
    return {{a, b}, dist(a,b)};
}

// se usa para el caso base, es O(1) en ese contexto
inline pair<ppt,fl> closest_brute(const vpt& P, int start, int end)
{
    pair<ppt,fl> p = {{P[start],P[start]}, INF};
    for (int i = start; i < end; ++i)
    for (int j = i+1; j < end; ++j)
        p = best_pair(p, mk_p(P[i], P[j]));
    return p;
}

// suena O(n^2), pero por matemagia divertida esto es realmente O(1) en el
// contexto de la función completa, y está por ahí de 7 y 12 la cota superior
inline pair<ppt,fl> strip_closest(const vpt& strip, pair<ppt,fl> closest)
{
    for (int i = 0; i < strip.size(); ++i)
    for (int j = i+1; j < strip.size() and (strip[j].y - strip[i].y) < closest.second; ++j)
        closest = best_pair(closest, mk_p(strip[i], strip[j]));
    return closest;
}

pair<ppt,fl> aux_closest_pair(const vpt& Px, const vpt& Py, int start, int end)
{
    // caso base O(1)
    if (end-start <= 3) return closest_brute(Px, start, end);

    int mid = start + ((end-start)>>1);
    pt midpt = Px[mid];

    // Dividimos en izq y der el ordenamiento vertical O(n)
    vpt Pyl, Pyr;
    Pyl.reserve(mid-start); Pyr.reserve(end-mid+1);
    for (const pt& p : Py)
        if (p.x <= midpt.x) Pyl.push_back(p);
        else Pyr.push_back(p);

    // divide
    auto closest = best_pair(aux_closest_pair(Px, Pyl, start, mid),
            aux_closest_pair(Px, Pyr, mid, end));

    // vencerás: la pareja es una de las mitades, o cruza la franja media
    vpt strip;
    for (const pt& p : Py) // O(n)
        if (ABS(p.x - midpt.x) < closest.second)
            strip.push_back(p);
    return strip_closest(strip, closest);
}

inline pair<ppt,fl> closest_pair(vpt& P) // O(nlogn) divide y vencerás
{
    vpt Px(P), Py(P);
    sort(ALL(Px), cmpx); sort(ALL(Py), cmpy);
    return aux_closest_pair(Px, Py, 0, P.size());
}

\end{minted}

	\needspace{2\baselineskip}
	{
		\phantomsection
		\subsection*{Complex Points}
		\addcontentsline{toc}{subsection}{Complex Points}
	}
	\begin{minted}{cpp}
#include <complex>
// definimos un punto como un complejo
//typedef long long C; // no nos sirve para algunos trucos como abs y arg
typedef long double C;
typedef complex<C> P;
#define X real()
#define Y imag()
#define C_EPS 1e-12

// no importa si C es punto flotante
inline bool operator<(P a, P b) { return a.X < b.X || (a.X == b.X && a.Y < b.Y); }
inline C cross(P a, P b) { return (conj(a)*b).Y; }
inline C dir(P u, P v, P w) { return cross(v-u,w-v); }
inline bool ccw(P u, P v, P w) { return dir(u,v,w) > C_EPS; }
inline bool cw(P u, P v, P w) { return dir(u,v,w) < -C_EPS; }
inline bool colin(P u, P v, P w)
{ C d = dir(u,v,w); return -C_EPS <= d && d <= C_EPS; }
int seg_intersect(P a, P b, P c, P d)
{ // 0.ab no toca cd.
  // 1.se tocan solo en un extremo.
  // 2.son colineales y se tocan infinitas veces.
  // 3.se tocan al interior de alguno exactamente una vez.
    if (colin(a,b,c) && colin(a,b,d))
    {
        if (c > d) swap(c,d);
        if (a > b) swap(a,b);
        if (abs(b-c) <= C_EPS || abs(d-a) <= C_EPS) return 1;
        if (a < c && c < b) return 2;
        if (c < a && a < d) return 2;
        return 0;
    }
    if (abs(a-c)<=C_EPS || abs(a-d)<=C_EPS || abs(b-c)<=C_EPS || abs(b-d)<=C_EPS)
        return 1;
    if (dir(a,b,c)*dir(a,b,d) < 0 && dir(c,d,a)*dir(c,d,b) < 0)
        return 3;
    return 0;
}
// solo si C es punto flotante
inline C dist(P a, P b) { return abs(b-a); }
inline C dist_pt_line(P l1, P l2, P p) { return cross(l1-p,l2-p)/dist(l1,l2); }
inline C angle_xaxis(P p) { return arg(p); }
inline P polr(C mag, C angle) { return polar(mag,angle); }
inline P rotated(P v, C angle) { return v*polar(1,angle); }

int main()
{
    P p = {3,-4};
    cout << "p: " << p.X << ' ' << p.Y << endl; // mostramos parte X y Y
    return 0;
}

\end{minted}

	\needspace{2\baselineskip}
	{
		\phantomsection
		\subsection*{Convex Hull}
		\addcontentsline{toc}{subsection}{Convex Hull}
	}
	\begin{minted}{cpp}
// algoritmo de graham
using ff = float;

// -1 reloj (derecho), +1 contrarreloj (izquierdo), 0 colineal
int giro(v2 a, v2 b, v2 c) {
    ff v = a.x*(b.y-c.y)+b.x*(c.y-a.y)+c.x*(a.y-b.y);
    if (v < 0) return -1;
    if (v > 0) return +1;
    return 0;
}

// regresa lista con el cierre convexo (CH)
vector<v2> cvxhull(vector<v2>& pts) {
    // caso trivial
    if (pts.size() <= 2) return pts;

    // punto mas chico en Y
    v2 p0 = *min_element(pts.begin(), pts.end(), [](const v2& a, const v2& b) {
        return make_pair(a.y, a.x) < make_pair(b.y, b.x);
    });
    // ord angular con respecto a p0
    sort(pts.begin(), pts.end(), [&p0](const v2& a, const v2& b){
        int g = giro(p0, a, b);
        if (g == 0) {
            // si son colineales, escoge primero el más cercano
            return (p0-a).sqrmag() < (p0-b).sqrmag();
        }
        return g > 0;
    });

    // ordenamos últimos colineales por distancia inversa para evitar que 'vaya y regrese'
    ll i = (ll)pts.size() - 1;
    while (i >= 0 && giro(p0, pts[i], pts.back()) == 0) i--;
    // invertimos los últimos puntos colineales para que aparezcan primero los lejanos
    reverse(pts.begin()+i+1, pts.end());

    vector<v2> cvx; // la pila con el CH
    for (ll j = 0; j < (ll)pts.size(); j++) {
        // eliminamos los que induzcan un giro inválido
        while (cvx.size() > 1 && giro(cvx[cvx.size()-2], cvx.back(), pts[j]) < 0)
            cvx.pop_back();
        cvx.push_back(pts[j]);
    }

    return cvx;
}

\end{minted}

	\needspace{2\baselineskip}
	{
		\phantomsection
		\subsection*{Delaunay}
		\addcontentsline{toc}{subsection}{Delaunay}
	}
	\begin{minted}{cpp}
//// Point.h

#define rep(i, a, b) for(int i = a; i < (b); ++i)
#define sz(x) (int)(x).size()

/**
 * Author: Ulf Lundstrom
 * Date: 2009-02-26
 * License: CC0
 * Source: My head with inspiration from tinyKACTL
 * Description: Class to handle points in the plane.
 * 	T can be e.g. double or long long. (Avoid int.)
 * Status: Works fine, used a lot
 */
#pragma once

//template <class T> int sgn(T x) { return (x > 0) - (x < 0); }
template<class T>
struct Point {
	typedef Point P;
	T x, y;
	explicit Point(T x=0, T y=0) : x(x), y(y) {}
	bool operator<(P p) const { return tie(x,y) < tie(p.x,p.y); }
	bool operator==(P p) const { return tie(x,y)==tie(p.x,p.y); }
	P operator+(P p) const { return P(x+p.x, y+p.y); }
	P operator-(P p) const { return P(x-p.x, y-p.y); }
	P operator*(T d) const { return P(x*d, y*d); }
	P operator/(T d) const { return P(x/d, y/d); }
	// T dot(P p) const { return x*p.x + y*p.y; }
	T cross(P p) const { return x*p.y - y*p.x; }
	T cross(P a, P b) const { return (a-*this).cross(b-*this); }
	T dist2() const { return x*x + y*y; }
	//double dist() const { return sqrt((double)dist2()); }
	// angle to x-axis in interval [-pi, pi]
	//double angle() const { return atan2(y, x); }
	//P unit() const { return *this/dist(); } // makes dist()=1
	//P perp() const { return P(-y, x); } // rotates +90 degrees
	//P normal() const { return perp().unit(); }
	// returns point rotated 'a' radians ccw around the origin
	//P rotate(double a) const {
		//return P(x*cos(a)-y*sin(a),x*sin(a)+y*cos(a)); }
	//friend ostream& operator<<(ostream& os, P p) {
		//return os << "(" << p.x << "," << p.y << ")"; }
};

//// Point.h


//// FastDelaunay.h

/**
 * Author: Philippe Legault
 * Date: 2016
 * License: MIT
 * Source: https://github.com/Bathlamos/delaunay-triangulation/
 * Description: Fast Delaunay triangulation.
 * Each circumcircle contains none of the input points.
 * There must be no duplicate points.
 * If all points are on a line, no triangles will be returned.
 * Should work for doubles as well, though there may be precision issues in 'circ'.
 * Returns triangles in order \{t[0][0], t[0][1], t[0][2], t[1][0], \dots\}, all counter-clockwise.
 * Time: O(n \log n)
 * Status: stress-tested
 */
#pragma once



typedef Point<ll> P;
typedef struct Quad* Q;
typedef __int128_t lll; // (can be ll if coords are < 2e4)
P arb(LLONG_MAX,LLONG_MAX); // not equal to any other point

struct Quad {
	Q rot, o; P p = arb; bool mark;
	P& F() { return r()->p; }
	Q& r() { return rot->rot; }
	Q prev() { return rot->o->rot; }
	Q next() { return r()->prev(); }
} *H;

bool circ(P p, P a, P b, P c) { // is p in the circumcircle?
	lll p2 = p.dist2(), A = a.dist2()-p2,
	    B = b.dist2()-p2, C = c.dist2()-p2;
	return p.cross(a,b)*C + p.cross(b,c)*A + p.cross(c,a)*B > 0;
}
Q makeEdge(P orig, P dest) {
	Q r = H ? H : new Quad{new Quad{new Quad{new Quad{0}}}};
	H = r->o; r->r()->r() = r;
	rep(i,0,4) r = r->rot, r->p = arb, r->o = i & 1 ? r : r->r();
	r->p = orig; r->F() = dest;
	return r;
}
void splice(Q a, Q b) {
	swap(a->o->rot->o, b->o->rot->o); swap(a->o, b->o);
}
Q connect(Q a, Q b) {
	Q q = makeEdge(a->F(), b->p);
	splice(q, a->next());
	splice(q->r(), b);
	return q;
}

pair<Q,Q> rec(const vector<P>& s) {
	if (sz(s) <= 3) {
		Q a = makeEdge(s[0], s[1]), b = makeEdge(s[1], s.back());
		if (sz(s) == 2) return { a, a->r() };
		splice(a->r(), b);
		auto side = s[0].cross(s[1], s[2]);
		Q c = side ? connect(b, a) : 0;
		return {side < 0 ? c->r() : a, side < 0 ? c : b->r() };
	}

#define H(e) e->F(), e->p
#define valid(e) (e->F().cross(H(base)) > 0)
	Q A, B, ra, rb;
	int half = sz(s) / 2;
	tie(ra, A) = rec({all(s) - half});
	tie(B, rb) = rec({sz(s) - half + all(s)});
	while ((B->p.cross(H(A)) < 0 && (A = A->next())) ||
	       (A->p.cross(H(B)) > 0 && (B = B->r()->o)));
	Q base = connect(B->r(), A);
	if (A->p == ra->p) ra = base->r();
	if (B->p == rb->p) rb = base;

#define DEL(e, init, dir) Q e = init->dir; if (valid(e)) \
		while (circ(e->dir->F(), H(base), e->F())) { \
			Q t = e->dir; \
			splice(e, e->prev()); \
			splice(e->r(), e->r()->prev()); \
			e->o = H; H = e; e = t; \
		}
	for (;;) {
		DEL(LC, base->r(), o);  DEL(RC, base, prev());
		if (!valid(LC) && !valid(RC)) break;
		if (!valid(LC) || (valid(RC) && circ(H(RC), H(LC))))
			base = connect(RC, base->r());
		else
			base = connect(base->r(), LC->r());
	}
	return { ra, rb };
}

vector<P> triangulate(vector<P> pts) {
	sort(all(pts));  assert(unique(all(pts)) == pts.end());
	if (sz(pts) < 2) return {};
	Q e = rec(pts).first;
	vector<Q> q = {e};
	int qi = 0;
	while (e->o->F().cross(e->F(), e->p) < 0) e = e->o;
#define ADD { Q c = e; do { c->mark = 1; pts.push_back(c->p); \
	q.push_back(c->r()); c = c->next(); } while (c != e); }
	ADD; pts.clear();
	while (qi < sz(q)) if (!(e = q[qi++])->mark) ADD;
	return pts;
}

//// FastDelaunay.h

\end{minted}

	\needspace{2\baselineskip}
	{
		\phantomsection
		\subsection*{KD Tree}
		\addcontentsline{toc}{subsection}{KD Tree}
	}
	\begin{minted}{cpp}
#include <bits/stdc++.h>
using namespace std;

using Coord = double; // tipo de coordenadas
constexpr Coord EPS = 1e-9;
const int MAX_D = 5; // numero de dimensiones máximo

struct KDTree {
    struct Point {
        Coord x[MAX_D];
        int id;
    };
    struct Node {
        Point p;
        int dim;
        Node *l, *r;
    };
    int D;
    Node* root;

    KDTree(int _D) : D(_D), root(nullptr) {}

    Node* build(vector<Point>& pts, int l, int r, int depth) {
        if (l >= r) return nullptr;
        int dim = depth % D;
        int m = (l + r) / 2;
        nth_element(pts.begin()+l, pts.begin()+m, pts.begin()+r,
            [dim](const Point& a, const Point& b){ return a.x[dim] < b.x[dim]; });
        Node* node = new Node{pts[m], dim, nullptr, nullptr};
        node->l = build(pts, l, m, depth+1);
        node->r = build(pts, m+1, r, depth+1);
        return node;
    }

    void build(vector<Point>& pts) { root = build(pts, 0, pts.size(), 0); }

    // ---- Distance function ----
    Coord dist(const Point& a, const Point& b) {
    #ifdef MANHATTAN
        Coord d = 0;
        for (int i=0;i<D;i++) d += abs(a.x[i]-b.x[i]);
        return d;
    #else
        Coord d = 0;
        for (int i=0;i<D;i++){ Coord diff=a.x[i]-b.x[i]; d += diff*diff; }
        return d;
    #endif
    }

    // ---- Nearest neighbor ----
    void nearest(Node* node, const Point& target, Point& best, Coord& bestd) {
        if (!node) return;
        Coord d = dist(node->p, target);
    #ifdef MANHATTAN
        if (d + EPS < bestd) { bestd = d; best = node->p; }
    #else
        if (d < bestd - EPS) { bestd = d; best = node->p; }
    #endif
        int dim = node->dim;
        Node *first = target.x[dim] < node->p.x[dim] ? node->l : node->r;
        Node *second = first == node->l ? node->r : node->l;
        nearest(first, target, best, bestd);
        Coord diff = abs(target.x[dim] - node->p.x[dim]);
    #ifdef MANHATTAN
        if (diff < bestd - EPS) nearest(second, target, best, bestd);
    #else
        if (diff*diff < bestd - EPS) nearest(second, target, best, bestd);
    #endif
    }

    Point nearest(const Point& target) {
        Point best; Coord bestd = numeric_limits<Coord>::max();
        nearest(root, target, best, bestd);
        return best;
    }

    void range_query(Node* node, const Coord low[], const Coord high[], vector<Point>& out) {
        if (!node) return;
        const auto& P = node->p;
        bool inside = true;
        for (int i=0; i<D; i++) {
            if (P.x[i] < low[i] - EPS || P.x[i] > high[i] + EPS)
            { inside = false; break; }
        }
        if (inside) out.push_back(P);

        int dim = node->dim;
        if (low[dim] - EPS <= P.x[dim]) range_query(node->l, low, high, out);
        if (high[dim] + EPS >= P.x[dim]) range_query(node->r, low, high, out);
    }

    vector<Point> range_query(const vector<Coord>& low, const vector<Coord>& high) {
        vector<Point> res;
        Coord lo[MAX_D], hi[MAX_D];
        for (int i=0;i<D;i++){ lo[i] = low[i]; hi[i] = high[i]; }
        range_query(root, lo, hi, res);
        return res;
    }
};

int main() {
    ios::sync_with_stdio(false);
    cin.tie(nullptr);

    int n, D; cin >> n >> D; // n - #pts, D - #dims
    vector<KDTree::Point> pts(n);
    for (int i=0;i<n;i++) {
        for (int j=0;j<D;j++) cin >> pts[i].x[j];
        pts[i].id = i;
    }
    KDTree tree(D);
    tree.build(pts);

    // Example nearest neighbor query
    KDTree::Point q;
    for (int j=0;j<D;j++) cin >> q.x[j];
    auto nearest_pt = tree.nearest(q);
    cout << "Nearest id: " << nearest_pt.id << "\n";

    // Example range query
    vector<Coord> low(D), high(D);
    for (int j=0;j<D;j++) cin >> low[j];
    for (int j=0;j<D;j++) cin >> high[j];
    auto ans = tree.range_query(low, high);
    for (auto &p : ans) {
        cout << p.id;
        for (int j=0;j<D;j++) cout << " " << p.x[j];
        cout << "\n";
    }
}

\end{minted}

	\needspace{2\baselineskip}
	{
		\phantomsection
		\subsection*{Li Chao Tree}
		\addcontentsline{toc}{subsection}{Li Chao Tree}
	}
	\begin{minted}{cpp}
// O(nlogC) con C el máximo |x_i|
// Para funciones que se intersectan a lo más una vez, nos da el min de
// todas ellas en tiempo logaritmico (por ejemplo para encontrar el CH de
// varias rectas)
typedef long long C; // esto se puede cambiar
typedef complex<C> P;
#define X real()
#define Y imag()

C dot(P a, P b) {
    return (conj(a) * b).X;
}

// rectas por ejemplo, pero puede ser otra fn
C f(P a,  C x) {
    return dot(a, {x, 1});
}

const int maxn = 2e5; // max número de funciones 

P line[4 * maxn]; // segment tree del arreglo de lineas

void add_line(P nw, int v = 1, int l = 0, int r = maxn) {
    int m = (l + r) / 2;
    bool lef = f(nw, l) < f(line[v], l);
    bool mid = f(nw, m) < f(line[v], m);
    if(mid) {
        swap(line[v], nw);
    }
    if(r - l == 1) {
        return;
    } else if(lef != mid) {
        add_line(nw, 2 * v, l, m);
    } else {
        add_line(nw, 2 * v + 1, m, r);
    }
}

C get(int x, int v = 1, int l = 0, int r = maxn) {
    int m = (l + r) / 2;
    if(r - l == 1) {
        return f(line[v], x);
    } else if(x < m) {
        return min(f(line[v], x), get(x, 2 * v, l, m));
    } else {
        return min(f(line[v], x), get(x, 2 * v + 1, m, r));
    }
}

// para agregar líneas: add_line({m, b})
// para obtener el mínimo del arreglo de líneas en x: get_line(x)

\end{minted}

	\needspace{2\baselineskip}
	{
		\phantomsection
		\subsection*{Line}
		\addcontentsline{toc}{subsection}{Line}
	}
	\begin{minted}{cpp}
// linea 2d

using ff = float;

// distancia linea-punto, donde la línea se define por a->b
ff dist_lin_v2(v2 p, v2 a, v2 b) {
    ff A = cross(b-a, p-a);
    return fabs(A)/((b-a).norm());
}

// si se intersectan dos líneas (p->r y q->s), y en out dónde
bool intersect_lin_lin(v2 p, v2 r, v2 q, v2 s, v2 &out){
  double rxs = cross(r, s);
  if (fabs(rxs) < 1e-12) 
    return false;
  double t = cross(q - p, s) / rxs;
  out = { p.x + t * r.x,
          p.y + t * r.y };
  return true;
}

\end{minted}

	\needspace{2\baselineskip}
	{
		\phantomsection
		\subsection*{Minimum Enclosing Circle}
		\addcontentsline{toc}{subsection}{Minimum Enclosing Circle}
	}
	\begin{minted}{cpp}
// Tienes que definir EPS

#define rep(i, a, b) for(int i = a; i < (b); ++i)
#define sz(x) (int)(x).size()

/// Points.h

/**
 * Author: Ulf Lundstrom
 * Date: 2009-02-26
 * License: CC0
 * Source: My head with inspiration from tinyKACTL
 * Description: Class to handle points in the plane.
 * 	T can be e.g. double or long long. (Avoid int.)
 * Status: Works fine, used a lot
 */
#pragma once

//template <class T> int sgn(T x) { return (x > 0) - (x < 0); }
template<class T>
struct Point {
	typedef Point P;
	T x, y;
	explicit Point(T x=0, T y=0) : x(x), y(y) {}
	bool operator<(P p) const { return tie(x,y) < tie(p.x,p.y); }
	bool operator==(P p) const { return tie(x,y)==tie(p.x,p.y); }
	P operator+(P p) const { return P(x+p.x, y+p.y); }
	P operator-(P p) const { return P(x-p.x, y-p.y); }
	P operator*(T d) const { return P(x*d, y*d); }
	P operator/(T d) const { return P(x/d, y/d); }
	//T dot(P p) const { return x*p.x + y*p.y; }
	T cross(P p) const { return x*p.y - y*p.x; }
	T cross(P a, P b) const { return (a-*this).cross(b-*this); }
	T dist2() const { return x*x + y*y; }
	double dist() const { return sqrt((double)dist2()); }
	// angle to x-axis in interval [-pi, pi]
	//double angle() const { return atan2(y, x); }
	//P unit() const { return *this/dist(); } // makes dist()=1
	P perp() const { return P(-y, x); } // rotates +90 degrees
	//P normal() const { return perp().unit(); }
	// returns point rotated 'a' radians ccw around the origin
	//P rotate(double a) const {
		//return P(x*cos(a)-y*sin(a),x*sin(a)+y*cos(a)); }
	//friend ostream& operator<<(ostream& os, P p) {
		//return os << "(" << p.x << "," << p.y << ")"; }
};

/// circumcircle.h

/**
 * Author: Ulf Lundstrom
 * Date: 2009-04-11
 * License: CC0
 * Source: http://en.wikipedia.org/wiki/Circumcircle
 * Description:\\
\begin{minipage}{75mm}
    The circumcirle of a triangle is the circle intersecting all three
    vertices. ccRadius returns the radius of the circle going through points A,
    B and C and ccCenter returns the center of the same circle.
\end{minipage}
\begin{minipage}{15mm}
\vspace{-2mm}
\includegraphics[width=\textwidth]{content/geometry/circumcircle}
\end{minipage}
 * Status: tested
 */
// Circumcírculo

//typedef Point<double> P;
//double ccRadius(const P& A, const P& B, const P& C) {
	//return (B-A).dist()*(C-B).dist()*(A-C).dist()/
			//abs((B-A).cross(C-A))/2;
//}

P ccCenter(const P& A, const P& B, const P& C) {
	P b = C-A, c = B-A;
	return A + (b*c.dist2()-c*b.dist2()).perp()/b.cross(c)/2;
}

// MinimumEnclosingCircle.h

/**
 * Author: Andrew He, chilli
 * Date: 2019-05-07
 * License: CC0
 * Source: folklore
 * Description: Computes the minimum circle that encloses a set of points.
 * Time: expected O(n)
 * Status: stress-tested
 */
// Este es el algoritmo de Welzl, no el de Megiddo (O(n) determinista)
pair<P, double> mec(vector<P> ps) {
	shuffle(all(ps), mt19937(time(0)));
	P o = ps[0];
	double r = 0, EPS = 1 + 1e-8;
	rep(i,0,sz(ps)) if ((o - ps[i]).dist() > r * EPS) {
		o = ps[i], r = 0;
		rep(j,0,i) if ((o - ps[j]).dist() > r * EPS) {
			o = (ps[i] + ps[j]) / 2;
			r = (o - ps[i]).dist();
			rep(k,0,j) if ((o - ps[k]).dist() > r * EPS) {
				o = ccCenter(ps[i], ps[j], ps[k]);
				r = (o - ps[i]).dist();
			}
		}
	}
	return {o, r};
}

\end{minted}

	\needspace{2\baselineskip}
	{
		\phantomsection
		\subsection*{Online Convex Hull}
		\addcontentsline{toc}{subsection}{Online Convex Hull}
	}
	\begin{minted}{cpp}
// O(logi) cada inserción, O(logi) consulta, O(nlogn) construcción total
struct pt
{
    ll x, y;
    pt(ll x_ = 0, ll y_ = 0) : x(x_), y(y_) {}
    bool operator<(const pt p) const // ord lexicográfico
    {
        if (x != p.x) return x < p.x;
        return y < p.y;
    }
    inline bool operator==(const pt p) const { return x == p.x and y == p.y; }
    inline pt operator+(const pt p) const { return pt(x+p.x, y+p.y); }
    inline pt operator-(const pt p) const { return pt(x-p.x, y-p.y); }
    inline pt operator*(const ll c) const { return pt(c*x, c*y); }
    inline ll operator*(const pt p) const { return x*p.x + y*p.y; } // producto interior
    inline int operator^(const pt p) const // signo del producto exterior
    {
        // NOTE: esto debe modificarse a según. estos productos fácilmente
        //       podrían ser muy grandes, por lo que se podría quizás adaptar
        //       como...
        //   --> long double l = x*(long double)p.y; // y análogamente para r
        //   --> long long l = x*p.y; // y análogamente para r
        //   --> dejarlo así y definir i128
        // x*p.y - y*p.x
        i128 l = mul_ll(x, p.y);
        i128 r = mul_ll(y, p.x);
        if (l == r) return 0;
        if (l < r) return -1;
        return 1;
    }
    friend istream& operator>>(istream& in, pt& p)
    {
        return in >> p.x >> p.y;
    }
};

bool ccw(pt p, pt q, pt r)
{
    return ((q-p)^(r-q)) > 0;
}

struct upper
{
    set<pt> se;
    set<pt>::iterator it;

    int is_under(pt p) // 1 -> in ; 2 -> border
    {
        it = se.lower_bound(p);
        if (it == se.end()) return 0;
        if (it == se.begin()) return p == *it ? 2 : 0;
        if (ccw(p, *it, *prev(it))) return 1;
        return ccw(p, *prev(it), *it) ? 0 : 2;
    }
    void insert(pt p)
    {
        if (is_under(p)) return;
        if (it != se.end()) while (next(it) != se.end() and !ccw(*next(it), *it, p))
            it = se.erase(it);
        if (it != se.begin()) while (--it != se.begin() and !ccw(p, *it, *prev(it)))
            it = se.erase(it);
        se.insert(p);
    }
};

struct dyn_hull
{
    upper U,L;

    int is_inside(pt p) // 1 -> in ; 2 -> border
    {
        int u = U.is_under(p), l = L.is_under({-p.x, -p.y});
        if (!u or !l) return 0;
        return max(u,l);
    }
    void insert(pt p)
    {
        U.insert(p);
        L.insert({-p.x, -p.y});
    }
    int size()
    {
        int ans = U.se.size() + L.se.size();
        return ans <= 2 ? ans/2 : ans-2;
    }
};


\end{minted}

	\needspace{2\baselineskip}
	{
		\phantomsection
		\subsection*{Picks Theorem}
		\addcontentsline{toc}{subsection}{Picks Theorem}
	}
	\begin{minted}{cpp}
// Si tenemos un polígono con coordenadas enteras y conocemos el número de
// puntos (con coordenadas enteras) internos y en la frontera, podemos obtener
// su área en O(1):
int pick_area(int internos, int en_frontera)
{
    return internos + en_frontera/2 - 1;
}

// Podemos obtener los puntos con coordenadas enteras en la frontera en O(n):
int boundary_points(const vector<pair<int, int>>& P) {
    int n = P.size();
    int B = 0;
    for (int i = 0; i < n; ++i) {
        int j = (i + 1) % n;
        int dx = abs(P[i].first - P[j].first);
        int dy = abs(P[i].second - P[j].second);
        B += __gcd(dx, dy); // tecnicamente O(log(min(dx,dy)))
    }
    return B;
}

// Si conocemos el área y los puntos en la frontera, podemos obtener internos.
// Tomamos 2 veces el área porque podemos obtenerla sin floats con shoelace
int pick_internos(int area2x, int en_frontera)
{
    return (area2x - en_frontera + 2)/2;
}

// asi...
ll shoelace_area2x(const vector<pair<int,int>>& P)
{
    int n = P.size();
    ll area2x = 0;
    for (int i = 0; i < n; ++i)
    {
        int j = (i+1)%n;
        area2x += 1LL * P[i].first * P[j].second;
        area2x -= 1LL * P[j].first * P[i].second;
    }
    return abs(area2x);
}

\end{minted}

	\needspace{2\baselineskip}
	{
		\phantomsection
		\subsection*{Points}
		\addcontentsline{toc}{subsection}{Points}
	}
	\begin{minted}{cpp}
// vector 2d

using ff = float;

struct v2 {
    ff x, y;
    v2(ff x, ff y): x(x), y(y) {}
    bool operator==(v2 const& t) const {
        return x == t.x && y == t.y;
    }
    v2& operator+=(const v2 &t) {
        x += t.x;
        y += t.y;
        return *this;
    }
    v2& operator-=(const v2 &t) {
        x -= t.x;
        y -= t.y;
        return *this;
    }
    v2& operator*=(ff t) {
        x *= t;
        y *= t;
        return *this;
    }
    v2& operator/=(ff t) {
        x /= t;
        y /= t;
        return *this;
    }
    v2 operator+(const v2 &t) const {
        return v2(*this) += t;
    }
    v2 operator-(const v2 &t) const {
        return v2(*this) -= t;
    }
    v2 operator*(ff t) const {
        return v2(*this) *= t;
    }
    v2 operator/(ff t) const {
        return v2(*this) /= t;
    }
    ff mag() const {
        return sqrt(x*x + y*y);
    }
    ff sqrmag() const {
        return x*x + y*y;
    }
};

v2 operator*(ff a, v2 b) {
    return b * a;
}

inline ff dot(v2 a, v2 b) {
    return a.x * b.x + a.y * b.y;
}

inline ff cross(v2 a, v2 b) {
    return a.x*b.y - a.y*b.x;
}

\end{minted}

	\needspace{2\baselineskip}
	{
		\phantomsection
		\subsection*{Shoelace Formula}
		\addcontentsline{toc}{subsection}{Shoelace Formula}
	}
	\begin{minted}{cpp}
// Obtiene en O(n) el área de un polígono simple con vértices en orden
// horario o anti-horario.
double shoelaceArea(const vector<pair<int, int>>& P) {
    int n = P.size();
    long long area = 0;
    for (int i = 0; i < n; ++i) {
        int j = (i + 1) % n;
        area += 1LL * P[i].first * P[j].second;
        area -= 1LL * P[j].first * P[i].second;
    }
    return abs(area) / 2.0;
}

\end{minted}

	\needspace{5\baselineskip}
	{
		\phantomsection
		\section*{Others}
		\markboth{OTHERS}{}
		\addcontentsline{toc}{section}{Others}
	}
	{
		\phantomsection
		\subsection*{All Submask}
		\addcontentsline{toc}{subsection}{All Submask}
	}
	\begin{minted}{cpp}
// O(3^n)
int dp[1 << 16];
int cost[1 << 16];
int main()
{
  int N = 16;
  for (int mask = 0; mask < N; mask++)
  {
    for (int submask = mask; submask; submask = (submask - 1) & mask)
    {
      dp[mask] = max(dp[mask], cost[submask] + dp[mask ^ submask]);
    }
  }
  return 0;
}
\end{minted}

	\needspace{2\baselineskip}
	{
		\phantomsection
		\subsection*{Binpow}
		\addcontentsline{toc}{subsection}{Binpow}
	}
	\begin{minted}{cpp}
ll binpow(ll a, ll b) {
  ll res = 1;
  while (b > 0) {
    if (b & 1)
      res = res * a;
    a = a * a;
    b >>= 1;
  }
  return res;
}
\end{minted}

	\needspace{2\baselineskip}
	{
		\phantomsection
		\subsection*{Coordinate Compression}
		\addcontentsline{toc}{subsection}{Coordinate Compression}
	}
	\begin{minted}{cpp}
vector<int> coordinate_compression(vector<int>& coordinates){
  vector<int> compression(coordinates.size()); 
  vector<int> aux = coordinates;
  sort(aux.begin(), aux.end());
  aux.resize(unique(aux.begin(), aux.end()) - aux.begin());
  for (int i = 0; i < n; ++i){
    compression[i] = lower_bound(aux.begin(), aux.end(), coordinates[i]) 
    - aux.begin();
  }
  return compression; 
}
\end{minted}

	\needspace{2\baselineskip}
	{
		\phantomsection
		\subsection*{Cover Points By Intervals}
		\addcontentsline{toc}{subsection}{Cover Points By Intervals}
	}
	\begin{minted}{cpp}
// Número de puntos x en 0 <= l <= r <= 1e18 tal que el 
// número de intervalos que lo cubren es igual a k, con k en [1...n]
void cover_points()
{
  int n; cin >> n;
  vector<pll> pairs(n); vector<ll> val;
  for (int i = 0; i < n; i++)
  {
    ll l, r;
    cin >> l >> r;
    pairs[i] = make_pair(l, r + 1);
    val.push_back(l);
    val.push_back(r + 1);
  }

  sort(all(val));
  val.resize(unique(all(val)) - val.begin());
  vector<ll> prefix(2 * n + 1);

  for (auto p : pairs)
  {
    int l = lower_bound(all(val), p.first) - val.begin();
    int r = lower_bound(all(val), p.second) - val.begin();
    prefix[l] += 1;
    prefix[r] -= 1;
  }
  for (int i = 1; i <= 2 * n; i++)
    prefix[i] += prefix[i - 1];
  
  vector<ll> res(n + 1);
  for (int i = 1; i < val.size(); i++)
  {
    ll num_points = val[i] - val[i - 1];
    res[prefix[i - 1]] += num_points;
  }
  // res[i] = El número de puntos que son cubiertos por i intervalos
}
\end{minted}

	\needspace{2\baselineskip}
	{
		\phantomsection
		\subsection*{FFT}
		\addcontentsline{toc}{subsection}{FFT}
	}
	\begin{minted}{cpp}
using C = double;
using CX = complex<C>;
const C PI = acos(-1);

// FFT + inverse FFT
void fft(vector<CX> & a, bool invert) {
    int n = a.size();

    for (int i = 1, j = 0; i < n; i++) {
        int bit = n >> 1;
        for (; j & bit; bit >>= 1)
            j ^= bit;
        j ^= bit;

        if (i < j)
            swap(a[i], a[j]);
    }

    for (int len = 2; len <= n; len <<= 1) {
        C ang = 2 * PI / len * (invert ? -1 : 1);
        CX wlen(cos(ang), sin(ang));
        for (int i = 0; i < n; i += len) {
            CX w(1);
            for (int j = 0; j < len / 2; j++) {
                CX u = a[i+j], v = a[i+j+len/2] * w;
                a[i+j] = u + v;
                a[i+j+len/2] = u - v;
                w *= wlen;
            }
        }
    }

    if (invert) {
        for (CX & x : a)
            x /= n;
    }
}

// multiplicación de polinomios en O(nlogn) [recibe coefs enteros]
vector<int> polyn_multiply(vector<int> const& a, vector<int> const& b) {
    vector<CX> fa(a.begin(), a.end()), fb(b.begin(), b.end());
    int n = 1;
    while (n < a.size() + b.size()) 
        n <<= 1;
    fa.resize(n);
    fb.resize(n);

    fft(fa, false);
    fft(fb, false);
    for (int i = 0; i < n; i++)
        fa[i] *= fb[i];
    fft(fa, true);

    vector<int> result(n);
    for (int i = 0; i < n; i++)
        result[i] = round(fa[i].real());
    return result;
}

// de cuántas maneras se puede escojer (i,j) tal que a[i]+b[j]=k?
int count_sum(vector<int> const& a, vector<int> const& b, int k)
{
    vector<int> pa = ..., pb = ...;
                       // aquí guardamos entradas de 0 (o min spg) hasta el
                       // valor max de a y b, donde en cada casilla arr[i]
                       // llenamos el número de ocurrencias de i.
    // si tenemos a=[1,2,3], b=[2,4], construimos
    //      pa = 0x^0 + 1x^1 + 1x^2 + 1x^3 + 0x^4
    //      pb = 0x^0 + 0x^1 + 1x^2 + 0x^3 + 1x^4
    // donde el coef. de x^i es el número de apariencias de i
    r = polyn_multiply(pa, pb);
    // el resultado es el coeficiente de x^k
    return r[k];
}

// dado vectores a y b de dimension n, encuentra el producto punto de a con
// cada shift circular de b.
vector<int> all_dot_shifts(vector<int> const& a, vector<int> const& b, int n)
{
    vector<int> fa(2*n), fb(2*n);
    fa = reverse(a) + [0]*n; // pseudocódigo pythonesco
    fb = b + b; // pseudocódigo pythonesco
    r = polyn_multiply(fa,fb);

    // así, r[k] nos da el producto punto de a con el shift (k-(n-1)) de b
}

// string matching con wildcards
vector<bool> str_match(const string& txt, const string& pat, char wildcard='*')
{
    int n = txt.size(), m = pat.size();
    int wilds = 0;
    for (char c : pat) if (c == wildcard) wilds++;
    vector<CX> a(n), b(m,0);
    for (int i = 0; i < n; i++)
    {
        C t = (2*PI*(txt[i]-'a'))/26;
        a[i] = {cos(t), sin(t)};
    }
    for (int i = 0; i < m; i++)
    {
        int j = m-i-1;
        if (pat[j] == wildcard) continue;
        C t = (2*PI*(pat[j]-'a'))/26;
        b[i] = {cos(t), sin(t)};
    }
    ab = polyn_multiply(a, b);
    vector<bool> r(n);
    for (int i = 0; i < n; i++) r[i] = ab[i] == (m - wilds);
    return r; // si r[i]=true entonces hay un match en la pos i
}

\end{minted}

	\needspace{2\baselineskip}
	{
		\phantomsection
		\subsection*{i128}
		\addcontentsline{toc}{subsection}{i128}
	}
	\begin{minted}{cpp}
// si no esta disponible
// typedef __int128_t lll;

// algunas operaciones que simulan int de 128 bits
using u64 = unsigned long long;

struct u128 { u64 hi, lo; };

// multiplica dos unsigned long long sin desborde
u128 mul_u64(u64 a, u64 b) {
    constexpr u64 MASK32 = 0xFFFFFFFF;

    u64 a_lo = a & MASK32;
    u64 a_hi = a >> 32;
    u64 b_lo = b & MASK32;
    u64 b_hi = b >> 32;

    u64 lo_lo = a_lo * b_lo;
    u64 hi_lo = a_hi * b_lo;
    u64 lo_hi = a_lo * b_hi;
    u64 hi_hi = a_hi * b_hi;

    u64 cross = hi_lo + lo_hi;
    bool carry1 = cross < hi_lo;

    u64 cross_lo = cross << 32;
    u64 lo = lo_lo + cross_lo;
    bool carry2 = lo < lo_lo;

    u64 hi = hi_hi + (cross >> 32) + carry1 + carry2;

    return { hi, lo };
}

// int de 128 bits con signo
struct i128
{
    u64 hi, lo;
    bool neg;
    inline bool operator==(const i128& o) const
    {
        return hi == o.hi and lo == o.lo and (neg == o.neg or (hi == 0 and lo == 0));
    }
    bool operator<(const i128& o) const
    {
        if (*this == o) return false;
        if (hi == 0 and lo == 0) return !o.neg;
        if (o.hi == 0 and o.lo == 0) return neg;
        if (neg != o.neg) return neg;
        if (hi < o.hi) return !neg;
        if (hi > o.hi) return neg;
        if (lo < o.lo) return !neg;
        if (lo > o.lo) return neg;
        return false;
    }
};

// multiplica dos long long sin desborde
i128 mul_ll(ll a, ll b)
{
    bool neg = (a<0) ^ (b<0);
    u128 res = mul_u64(llabs(a), llabs(b));
    return { res.hi, res.lo, neg };
}

\end{minted}

	\needspace{2\baselineskip}
	{
		\phantomsection
		\subsection*{Kadanes}
		\addcontentsline{toc}{subsection}{Kadanes}
	}
	\begin{minted}{cpp}
// Given an array of numbers a[0...n] find the subarray 
// a[l...r] with the maximum sum 
int ans = a[0], ans_l = 0, ans_r = 0;
int sum = 0, minus_pos = -1;
for (int r = 0; r < n; ++r) {
  sum += a[r];
  if (sum > ans) {
    ans = sum;
    ans_l = minus_pos + 1;
    ans_r = r;
  }
  if (sum < 0) {
    sum = 0;
    minus_pos = r;
  }
}
\end{minted}

	\needspace{2\baselineskip}
	{
		\phantomsection
		\subsection*{Matrix}
		\addcontentsline{toc}{subsection}{Matrix}
	}
	\begin{minted}{cpp}

template <typename T>
void matmul(vector<vector<T>> &a, vector<vector<T>> b)
{
    int n = a.size(), m = a[0].size(), p = b[0].size();
    assert(m == b.size());
    vector<vector<T>> c(n, vector<T>(p));
    for (int i = 0; i < n; i++)
        for (int j = 0; j < p; j++)
            for (int k = 0; k < m; k++)
                // Mult with mod define
                c[i][j] = (c[i][j] + a[i][k] * b[k][j]) % MOD;

    a = c;
}

template <typename T>
struct Matrix
{
    vector<vector<T>> mat;
    Matrix() {}
    Matrix(vector<vector<T>> a) { mat = a; }
    Matrix(int n, int m)
    {
        mat.resize(n);
        for (int i = 0; i < n; i++)
            mat[i].resize(m);
    }
    int rows() const { return mat.size(); }
    int cols() const { return mat[0].size(); }

    vector<T> &operator[](int i) { return mat[i]; }
    const vector<T> &operator[](int i) const { return mat[i]; }

    void makeiden()
    {
        for (int i = 0; i < rows(); i++)
            mat[i][i] = 1;
    }
    Matrix operator*=(const Matrix &b)
    {
        matmul(mat, b.mat);
        return *this;
    }
    Matrix operator*(const Matrix &b) { return Matrix(*this) *= b; }
    void print() const
    {
        for (int i = 0; i < rows(); i++)
        {
            for (int j = 0; j < cols(); j++)
            {
                cout << mat[i][j] << ' ';
            }
            cout << '\n';
        }
    }
};

\end{minted}

	\needspace{2\baselineskip}
	{
		\phantomsection
		\subsection*{Matrix Exponentiation}
		\addcontentsline{toc}{subsection}{Matrix Exponentiation}
	}
	\begin{minted}{cpp}

// Uso: Matrix<T> M = matPow<T>(A, n - 1);
template <typename T>
Matrix<T> matPow(Matrix<T> &base, ll p)
{
  Matrix<T> ans(base.rows(), base.cols());
  ans.makeiden();
  while (p)
  {
    if (p & 1)
      ans *= base;

    base *= base;
    p >>= 1;
  }
  return ans;
}
\end{minted}

	\needspace{2\baselineskip}
	{
		\phantomsection
		\subsection*{Matrix with array}
		\addcontentsline{toc}{subsection}{Matrix with array}
	}
	\begin{minted}{cpp}
#include <iostream>
#include <algorithm>
using namespace std;

const long long MOD = 1000000007;
// Multiplicación de matrices 2x2: C = A * B
void mult2x2(const long long A[2][2], const long long B[2][2], long long C[2][2]) {
    long long temp00 = (A[0][0] * B[0][0] + A[0][1] * B[1][0]) % MOD;
    long long temp01 = (A[0][0] * B[0][1] + A[0][1] * B[1][1]) % MOD;
    long long temp10 = (A[1][0] * B[0][0] + A[1][1] * B[1][0]) % MOD;
    long long temp11 = (A[1][0] * B[0][1] + A[1][1] * B[1][1]) % MOD;
    C[0][0] = temp00;
    C[0][1] = temp01;
    C[1][0] = temp10;
    C[1][1] = temp11;
}
// Exponenciación binaria de la matriz 2x2 M a la potencia exp. El resultado se almacena en 'result'.
void matPow2x2(const long long M[2][2], long long exp, long long result[2][2]) {
    // Inicializamos 'result' como la identidad.
    result[0][0] = 1; result[0][1] = 0;
    result[1][0] = 0; result[1][1] = 1;
    
    long long base[2][2] = { {M[0][0], M[0][1]}, {M[1][0], M[1][1]} };
    
    while(exp > 0) {
        if(exp & 1) {
            long long temp[2][2];
            mult2x2(result, base, temp);
            result[0][0] = temp[0][0]; result[0][1] = temp[0][1];
            result[1][0] = temp[1][0]; result[1][1] = temp[1][1];
        }
        long long temp[2][2];
        mult2x2(base, base, temp);
        base[0][0] = temp[0][0]; base[0][1] = temp[0][1];
        base[1][0] = temp[1][0]; base[1][1] = temp[1][1];
        exp >>= 1;
    }
}

// s_0 y s_1 son los casos base y n es el número de pasos a avanzar a partir de s1.
long long fib(long long n, long long s_0, long long s_1) {
    // Si n es 0, la matriz identidad deja la columna [s_1, s_0] y se devuelve s_1.
    if(n == 0) return s_1 % MOD;
    long long M[2][2] = { {1, 1}, {1, 0} };
    long long power[2][2];
    matPow2x2(M, n, power);
    // Multiplicamos la matriz resultante por la columna [s_1, s_0]:
    long long res = (power[0][0] * s_1 + power[0][1] * s_0) % MOD;
    return res;
}




\end{minted}

	\needspace{2\baselineskip}
	{
		\phantomsection
		\subsection*{Modular binpow}
		\addcontentsline{toc}{subsection}{Modular binpow}
	}
	\begin{minted}{cpp}
ll modpow(ll a, ll b, ll m) {
  a %= m;
  ll res = 1;
  while (b > 0) {
    if (b & 1)
      res = res * a % m;
    a = a * a % m;
    b >>= 1;
  }
  return res;
}
\end{minted}

	\needspace{2\baselineskip}
	{
		\phantomsection
		\subsection*{Preifix Sum 2D}
		\addcontentsline{toc}{subsection}{Preifix Sum 2D}
	}
	\begin{minted}{cpp}
// Se le suma uno MAX_SIDE para trabajar con 1-indexed y
// que la primera fila y columnas estén llenas de 0 y no haya problema
// al acceder a incides anteriores
constexpr int MAX_SIDE = 1000;
int mtxPrefix[MAX_SIDE + 1][MAX_SIDE + 1];
int mtx[MAX_SIDE + 1][MAX_SIDE + 1];
int main()
{
  int N, Q;
  cin >> N >> Q;

  for(int i = 1; i<= N; i++)
    for(int j = 1; j <= N; j++)
      cin >> mtx[i][j]; 
  
  // Construye la prefix matrix
  for (int i = 1; i <= N; i++)
    for (int j = 1; j <= N; j++)
      mtxPrefix[i][j] = mtx[i][j] + mtxPrefix[i - 1][j] 
                        + mtxPrefix[i][j - 1] - mtxPrefix[i - 1][j - 1];

  // Responder a querys donde te piden la suma de la matriz con esquina
  // superior izquierda (from_row, from_col) y esquina inferior derecha
  // (to_row, to_col), en este caso 1-indexed
  while (Q--)
  {
    int from_row, to_row, from_col, to_col;
    cin >> from_row >> from_col >> to_row >> to_col;
    ll total_sum = mtxPrefix[to_row][to_col] - mtxPrefix[from_row - 1][to_col] -
                   mtxPrefix[to_row][from_col - 1] +
                   mtxPrefix[from_row - 1][from_col - 1];
    cout << total_sum << endl;
  }
}
\end{minted}

	\needspace{2\baselineskip}
	{
		\phantomsection
		\subsection*{Static Range Minimum Query}
		\addcontentsline{toc}{subsection}{Static Range Minimum Query}
	}
	\begin{minted}{cpp}
// O(n) construccion, O(1) query
template<typename T> struct rmq
{
    vector<T> v;
    ll n; static const int b = 62; // cambia a 30 si usamos int en lugar de ll
    vector<ll> mask, t;

    // la operación se puede cambiar si el resultado es elemento del cjto
    ll op(ll x, ll y) { return v[x] < v[y] ? x : y; }
    ll msb(ll x) { return __builtin_clz(1)-__builtin_clz(x); }
    ll small (ll r, ll sz = b) { return r-msb(mask[r]&((1<<sz)-1)); }
    rmq(const vector<T>& v_) : v(v_), n(v.size()), mask(n), t(n)
    {
        for (ll i = 0, at = 0; i < n; mask[i++] = at |= 1)
        {
            at = (at<<1)&((1ll<<b)-1);
            while (at and op(i, i-msb(at&-at)) == i)
                at ^= at&-at;
        }
        for (ll i = 0; i < n/b; i++)
            t[i] = small(b*i+b-1);
        for (ll j = 1; (1ll<<j) <= n/b; j++)
        for (ll i = 0; i+(1ll<<j) <= n/b; i++)
                t[n/b*j+1] = op(t[n/b*(j-1)+i], t[n/b*(j-1)+i+(1ll<<(j-1))]);
    }
    T query(ll l, ll r)
    {
        if (r-l+1 <= b) return v[small(r,r-l+1)];
        ll ans = op(small(l+b-1), small(r));
        ll x = l/b+1, y = r/b-1;
        if (x <= y)
        {
            ll j = msb(y-x+1);
            ans = op(ans, op(t[n/b*j+x], t[n/b*j+y-(1ll<<j)+1]));
        }
        return v[ans];
    }
};

\end{minted}

	\needspace{2\baselineskip}
	{
		\phantomsection
		\subsection*{Ternary search}
		\addcontentsline{toc}{subsection}{Ternary search}
	}
	\begin{minted}{cpp}
#include <iostream>
#include <cmath>
#include <iomanip>
using namespace std;
const double PI = 3.14159265358979323846;

double f(double x) {
    return -(x - 2) * (x - 2) + PI; 
}

double busqueda_ternaria_max(double l, double r, double epsilon = 1e-7) {
    while (r - l > epsilon) {
        double tercio1 = l + (r - l) / 3;
        double tercio2 = r - (r - l) / 3;
        if (f(tercio1) < f(tercio2))
            l = tercio1;  // El máximo está en [tercio1, r]
        else
            r = tercio2;  // El máximo está en [l, tercio2]
    }
    return (l + r) / 2; // Aproximación del máximo
}

int main() {
    double l = 0, r = 4;
    double max_x = busqueda_ternaria_max(l, r);
    cout << fixed << setprecision(7);
    cout << "Máximo en x ≈ " << max_x << ", f(x) ≈ " << f(max_x) << endl;
    return 0;
}

\end{minted}

\end{multicols*}

\end{document}
